\documentclass[leqno]{article}
% Shared QBT/Soln pack preamble for resources.danieljsmith.org.
\usepackage{graphicx}
\usepackage{dsfont}
\usepackage{mathtools}
\usepackage{amssymb}
\usepackage{amsmath}
\usepackage{mleftright}
\usepackage{amsthm}
\usepackage{float}
\usepackage{ragged2e}
\usepackage{tensor}
\usepackage{fancyhdr}
\usepackage{empheq}
\usepackage[most]{tcolorbox}
\usepackage{enumitem}
\usepackage{dirtytalk}
\usepackage{marginnote}
\usepackage{microtype}
\usepackage[
  a4paper,
  left=0.8in,
  right=1in,
  top=1in,
  bottom=0.5in,
  headheight=14pt,
  headsep=0.4in
]{geometry}
\setlength{\parindent}{0cm}

% TikZ / PGFPlots
\usepackage{pgfplots}
\pgfplotsset{compat=1.18}
\usepgfplotslibrary{fillbetween, polar}
\usetikzlibrary{calc, positioning}
\usetikzlibrary{angles, quotes}
\usetikzlibrary{arrows.meta}
\usetikzlibrary{decorations.pathreplacing, decorations.markings}
\usetikzlibrary{patterns}
\usetikzlibrary{intersections}

% Links and contents. The contents follow the TMUA papers: compact entries,
% no dot leaders, and each question's mark total in grey after its label.
\usepackage{tocloft}
\tocloftpagestyle{fancy}
\setlength{\cftbeforesubsecskip}{2pt}
\setlength{\cftsubsecindent}{1.5em}
\setlength{\cftsubsecnumwidth}{0pt}
\renewcommand{\cftsubsecfont}{\small\color{black}}
\renewcommand{\cftsubsecpagefont}{\small\color{black}}
\renewcommand{\cftsubsecleader}{\hfill}
\usepackage[colorlinks=true,linkcolor=black,citecolor=magenta,urlcolor=black]{hyperref}%
\urlstyle{same}

% Pack macros
% \marks also adds its value to the current question's total in the .aux
% file, so the contents show each total from the second pass. The write
% happens at shipout, so a box typeset more than once counts once. The
% pack's mark total for the contents-page title is summed the same way.
\makeatletter
\newcount\djs@totalmarks
\newcommand{\djsAddMarks}[2]{%
  \global\advance\djs@totalmarks#2\relax
  \@ifundefined{djs@marks@#1}%
    {\expandafter\xdef\csname djs@marks@#1\endcsname{#2}}%
    {\expandafter\xdef\csname djs@marks@#1\endcsname
      {\the\numexpr\csname djs@marks@#1\endcsname+#2\relax}}}
\RenewDocumentCommand{\marks}{o m}{%
  \protected@write\@auxout{}{\string\djsAddMarks{\number\value{enumi}}{#2}}%
  \IfValueTF{#1}%
    {\marginnote{\raggedright\textbf{[#2]}}[#1]}%
    {\marginnote{\raggedright\textbf{[#2]}}}%
}
\makeatother

\newcommand{\vspacerule}[2]{%
  \par\vspace*{#1}%
  \noindent\hrulefill\par
  \vspace*{#2}%
}

\definecolor{scarlet}{RGB}{205, 40, 75}
% Site colours for the contents-page title: the accent, the rule colour, a
% stronger rule colour for the path separator, and the muted text colour.
\definecolor{djsBurgundy}{RGB}{118, 59, 67}
\definecolor{djsHairline}{RGB}{216, 213, 205}
\definecolor{djsSlash}{RGB}{170, 166, 158}
\definecolor{djsMuted}{RGB}{96, 101, 106}

% Worked-solution box, adapted from the TMUA worked-solutions box. A white breakable box with a left scarlet rule and no surrounding
% frame or tint. A part that breaks fills the rest of its page, so the rule
% runs to the bottom margin; the unbroken or last part ends in a short
% horizontal foot rule so a solution visibly stops. Unlike the TMUA box there
% is no answer label. The box does not grow to the left, so the rule sits at
% the text edge; growing to the right keeps the usable line width close to the
% old tcolorbox Solution box.
% Spacing inside the box: paragraphs are 6pt apart, so source blank lines show
% (solutions put one after each display and none between a lead-in and its
% display). Every display gets 6pt above and 10pt below, whatever the length
% of its lead-in line; align rows are spread by an 8pt \jot; solution parts
% are 14pt apart and sub-parts 8pt; \left and \right space like ordinary
% brackets (mleftright); no single line is stranded at a page break; and
% neither the Solution label nor the line before a list ends a page. No
% skip can shrink, so full pages are not squeezed.
\newtcolorbox{djsSolution}{%
  enhanced, breakable, frame hidden, sharp corners, colback=white,
  height fixed for=first and middle, height fill,
  extras unbroken and last={natural height},
  borderline west={2.2pt}{0pt}{scarlet},
  left=11pt, right=0pt, top=5pt, bottom=13pt,
  before skip=13pt, after skip=4pt,
  grow to right by=10mm,
  overlay unbroken and last={%
    \fill[scarlet] (frame.south west)
      rectangle ([xshift=3.5cm, yshift=2.2pt]frame.south west);},
  before upper={{\color{scarlet}\bfseries Solution}\par\nobreak\vspace{4pt}%
    \setlength{\parskip}{6pt plus 2pt}%
    \setlength{\abovedisplayskip}{6pt plus 2pt}%
    \setlength{\abovedisplayshortskip}{6pt plus 2pt}%
    \setlength{\belowdisplayskip}{10pt plus 2pt}%
    \setlength{\belowdisplayshortskip}{10pt plus 2pt}%
    \setlength{\jot}{8pt}%
    \setlist{parsep=6pt plus 2pt, topsep=6pt plus 2pt, beginpenalty=10000}%
    \setlist[enumerate,1]{itemsep=14pt plus 2pt}%
    \setlist[enumerate,2]{itemsep=8pt plus 2pt}%
    \mleftright
    \clubpenalty=10000 \widowpenalty=10000 \displaywidowpenalty=10000\relax}}

% Optional smaller figure inside a solution box. \scalebox shrinks node text
% with the drawing. Not applied to existing figures.
\newcommand{\SolnFigureScale}{0.85}
\newsavebox{\djsSolnFigureBox}
\newenvironment{djsSolutionFigure}
  {\par\vspace{4pt}\begin{center}\begin{lrbox}{\djsSolnFigureBox}}
  {\end{lrbox}\scalebox{\SolnFigureScale}{\usebox{\djsSolnFigureBox}}%
   \end{center}\vspace{2pt}\par}

\newcommand{\cxl}[2]{%
  \tikz[baseline=(X.base)]{%
    \node[inner sep=0pt, outer sep=0pt, text=#1] (X) {$#2$};
    \draw[#1, line width=0.4pt] (X.south west)--(X.north east);
  }%
}

% Page-style hook run at the contents page and each question's opening page.
% A no-op for QBT packs; \djsSolnHeader points it at the fancy header.
\newcommand{\djsQuestionPageStyle}{}

% A contents entry: "Question N" in a box wide enough for two digits, so the
% mark totals line up, then the total in grey. Bookmarks keep "Question N".
\makeatletter
\newlength{\djsTocLabelWidth}
\newcommand{\djsTocMarks}[1]{%
  \@ifundefined{djs@marks@#1}{}{%
    {\footnotesize\itshape\color{black!45}%
     \csname djs@marks@#1\endcsname~mark%
     \ifnum\csname djs@marks@#1\endcsname=1 \else s\fi}}}
\newcommand{\djsTocEntry}[1]{%
  \texorpdfstring
    {\settowidth{\djsTocLabelWidth}{Question 00}%
     \makebox[\dimexpr\djsTocLabelWidth+0.9em\relax][l]{Question #1}%
     \djsTocMarks{#1}}%
    {Question #1}}
\makeatother

% The question count for the contents-page title is written to the .aux at
% the end of the document, so it also appears from the second pass.
\newcounter{djsQuestions}
\newcommand{\questionitem}{%
  \djsQuestionPageStyle
  \item\phantomsection
  \stepcounter{djsQuestions}%
  \addcontentsline{toc}{subsection}{\protect\djsTocEntry{\number\value{enumi}}}%
}
\makeatletter
\newcommand{\djsSetQuestionCount}[1]{\gdef\djs@questioncount{#1}}
\AtEndDocument{\immediate\write\@auxout
  {\string\djsSetQuestionCount{\number\value{djsQuestions}}}}
\makeatother

\renewcommand{\contentsname}{Questions}

% \djsPackHeader[<link content tag>]{<left header>}; the tag says which
% document a click on the header link came from.
\newcommand{\djsPackHeader}[2][questions]{%
  \pagestyle{fancy}%
  \fancyhead{}%
  \fancyfoot{}%
  \renewcommand{\headrulewidth}{0.9pt}%
  \lhead{#2}%
  \chead{}%
  \rhead{\href{https://danieljsmith.org/?utm_source=pdf&utm_campaign=further-maths&utm_content=#1}{danieljsmith.org}}%
}

% The header's topic is also the contents-page title, under which a label
% says which booklet this is.
\newcommand{\djsPackTitle}{}
\newcommand{\djsBookletLabel}{}
\newcommand{\djsQbtHeader}[1]{%
  \djsPackHeader{#1}%
  \renewcommand{\djsPackTitle}{#1}%
  \renewcommand{\djsBookletLabel}{Questions by topic}%
}

% The colophon at the foot of page one, independent of the page style: the
% contents-page title's hairline led by a short scarlet rule, then the date on
% the left and a link to the site on the right, all inside the bottom margin.
% Its sizes are fixed, so it matches the TMUA papers' colophon. The document
% supplies the date; recompiling never changes it.
\makeatletter
\newcommand{\djsLastUpdated}[1]{%
  \AtBeginDocument{\AddToHookNext{shipout/foreground}{\djs@colophon{#1}}}}
\newcommand{\djs@colophon}[1]{%
  \put(0.8in,-\dimexpr\paperheight-0.47in\relax){%
    \makebox[0pt][l]{\color{djsHairline}%
      \rule[0.9pt]{\dimexpr\paperwidth-1.8in\relax}{0.6pt}}%
    \makebox[0pt][l]{\color{scarlet}\rule{3.5cm}{2.2pt}}}%
  \put(0.8in,-\dimexpr\paperheight-0.32in\relax){\makebox[0pt][l]{%
    \fontsize{9}{11}\selectfont Last updated #1}}%
  \put(\dimexpr\paperwidth-1in\relax,-\dimexpr\paperheight-0.32in\relax){%
    \makebox[0pt][r]{\fontsize{7}{8.5}\selectfont\color{djsMuted}%
      Latest version at \href{https://resources.danieljsmith.org/}%
        {\color{black}resources.danieljsmith.org}}}}
\makeatother

% Solutions packs show the header only on the contents page and each
% question's opening page, so it never appears part-way through a solution.
% \thispagestyle marks the next page shipped out, which is the question's
% page because every solution question starts after a \newpage.
\newcommand{\djsSolnHeader}[1]{%
  \djsPackHeader[solutions]{\textit{Solutions} - #1}%
  \pagestyle{empty}%
  \renewcommand{\djsQuestionPageStyle}{\thispagestyle{fancy}}%
  \renewcommand{\djsPackTitle}{#1}%
  \renewcommand{\djsBookletLabel}{Worked solutions}%
}

% \djsSitePath{<strand>}{<topic>}, called once after the header, names where
% the pack sits on the site; the contents-page title prints it above the pack
% title, e.g. CORE PURE / FURTHER CALCULUS. Without it the line is left out.
\makeatletter
\newif\ifdjs@sitepath
\newcommand{\djs@strand}{}
\newcommand{\djs@section}{}
\newcommand{\djsSitePath}[2]{%
  \djs@sitepathtrue
  \renewcommand{\djs@strand}{#1}%
  \renewcommand{\djs@section}{#2}%
}
\makeatother

% Further Mechanics packs call \djsMechanicsPack in their preamble, so that
% \djsFrontMatter prints the default value of g above the contents. The box
% keeps its own space above it (before skip is the whole gap below the title,
% not an addition to it). Without the box, \djsGravityNote is just extra space
% before the contents heading, bringing the gap above it close to the gap
% between the box and the heading on Further Mechanics pages.
\newcommand{\djsGravityNote}{\vspace{14pt}}
\newcommand{\djsMechanicsPack}{%
  \renewcommand{\djsGravityNote}{%
    \begin{tcolorbox}[sharp corners, colback=white, colframe=black,
      boxrule=1pt, left=8pt, right=8pt, top=6pt, bottom=6pt,
      before skip=20pt]
    {\bfseries Instructions}\par\vspace{4pt}
    Unless a question states otherwise, take $g = 9.8\,\mathrm{m\,s^{-2}}$
    wherever a numerical value of $g$ is required. Solutions also give final
    answers for $g = 10\,\mathrm{m\,s^{-2}}$ where these differ.
    \end{tcolorbox}%
    \vspace{16pt}}%
}

% Contents-page title: the site path in spaced burgundy capitals, the pack
% title, a hairline led by a short scarlet rule (the solution box's foot rule),
% then the booklet label with the question and mark counts. The counts come
% from the .aux, so they appear from the second pass.
\makeatletter
\newsavebox{\djs@eyebrowbox}
\newcommand{\djs@caps}[1]{\textls[180]{\MakeUppercase{#1}}}
\newcommand{\djs@eyebrow}{%
  \color{djsBurgundy}\footnotesize\bfseries
  \djs@caps{\djs@strand}%
  {\color{djsSlash}\hspace{0.55em}/\hspace{0.55em}}%
  \djs@caps{\djs@section}}
\newcommand{\djs@dot}{\hspace{0.45em}\textperiodcentered\hspace{0.45em}}
\newcommand{\djs@counts}{%
  \@ifundefined{djs@questioncount}{}{%
    \djs@questioncount~question\ifnum\djs@questioncount=1 \else s\fi
    \ifnum\djs@totalmarks>0
      \djs@dot\the\djs@totalmarks~mark\ifnum\djs@totalmarks=1 \else s\fi
    \fi}}
\newcommand{\djsContentsTitle}{%
  \ifdjs@sitepath
    \sbox\djs@eyebrowbox{\djs@eyebrow}%
    \ifdim\wd\djs@eyebrowbox>\linewidth
      \@latex@warning@no@line{The \string\djsSitePath\space line is wider
        than the text, so it wraps. Shorten its topic}%
    \fi
    {\raggedright\djs@eyebrow\par}%
    \vspace{7pt}%
  \fi
  {\raggedright\fontsize{24.88}{28}\selectfont\djsPackTitle\par}%
  \vspace{9pt}%
  \noindent\makebox[0pt][l]{\color{djsHairline}\rule[0.9pt]{\linewidth}{0.6pt}}%
  {\color{scarlet}\rule{3.5cm}{2.2pt}}\par
  \vspace{5pt}%
  {\small\color{djsMuted}\djsBookletLabel\hfill\djs@counts\par}%
  \vspace{8pt}%
}
\makeatother

\newcommand{\djsFrontMatter}{%
  \djsQuestionPageStyle
  \vspace*{6pt}%
  \djsContentsTitle
  \djsGravityNote
  \tableofcontents
  \newpage
}


\djsSolnHeader{Eigenvalues, Eigenvectors and Diagonalisation}
\djsSitePath{Core Pure}{Vectors, Matrices \& Linear Transformations}
\djsLastUpdated{27 October 2026}

\begin{document}
\djsFrontMatter
\begin{enumerate}[
  label=\textbf{\arabic*.},
  leftmargin=*,
  topsep=0pt,
  partopsep=0pt,
  parsep=0pt
]

% ---- Question 1 ----
\questionitem
% [Source: _QBT___Solns__Eigenvalues__Eigenvectors_and_Diagonalisation, Q1]

The matrix $\mathbf{A}$ is given by
\[\mathbf{A} = \begin{pmatrix} 3 & k \\ 2 & -2 \end{pmatrix}\]
where $k$ is a constant. \\

Given that $4$ is an eigenvalue of $\mathbf{A}$, find
\begin{itemize}
  \item the value of $k$,
  \item the other eigenvalue of $\mathbf{A}$,
  \item an eigenvector associated with each eigenvalue of $\mathbf{A}$.
\end{itemize}
\marks[-10pt]{6}

\vspace*{10pt}

\begin{djsSolution}
The eigenvalues satisfy the characteristic equation \(\det(\mathbf{A}-\lambda\mathbf{I})=0\), which is
\begin{align*}
\begin{vmatrix} 3-\lambda & k \\ 2 & -2-\lambda \end{vmatrix} &= 0\\
(3-\lambda)(-2-\lambda)-2k &= 0\\
\lambda^2-\lambda-6-2k &= 0
\end{align*}

Since \(4\) is an eigenvalue, substitute \(\lambda=4\):
\[
16-4-6-2k=0\quad\Longrightarrow\quad 6-2k=0\quad\Longrightarrow\quad k=3
\]

With \(k=3\) the characteristic equation becomes
\begin{align*}
\lambda^2-\lambda-12&=0\\
(\lambda-4)(\lambda+3)&=0
\end{align*}

so the other eigenvalue is \(\lambda=-3\).

For \(\lambda=4\), solve \((\mathbf{A}-4\mathbf{I})\begin{pmatrix}x\\y\end{pmatrix}=\mathbf{0}\):
\[
\begin{pmatrix} -1 & 3 \\ 2 & -6 \end{pmatrix}\begin{pmatrix}x\\y\end{pmatrix}=\begin{pmatrix}0\\0\end{pmatrix}
\]

Both rows give \(-x+3y=0\), so \(x=3y\). Taking \(y=1\), an eigenvector is
\[
\begin{pmatrix}3\\1\end{pmatrix}
\]

For \(\lambda=-3\), solve \((\mathbf{A}+3\mathbf{I})\begin{pmatrix}x\\y\end{pmatrix}=\mathbf{0}\):
\[
\begin{pmatrix} 6 & 3 \\ 2 & 1 \end{pmatrix}\begin{pmatrix}x\\y\end{pmatrix}=\begin{pmatrix}0\\0\end{pmatrix}
\]

Both rows give \(2x+y=0\), so \(y=-2x\). Taking \(x=1\), an eigenvector is
\[
\begin{pmatrix}1\\-2\end{pmatrix}
\]

So \(k=3\), the other eigenvalue is \(-3\), and eigenvectors are \(\begin{pmatrix}3\\1\end{pmatrix}\) for \(\lambda=4\) and \(\begin{pmatrix}1\\-2\end{pmatrix}\) for \(\lambda=-3\) (or any non-zero multiples of these).
\end{djsSolution}

\newpage

% ---- Question 2 ----
\questionitem
% [Source: _QBT___Solns__Eigenvalues__Eigenvectors_and_Diagonalisation, Q2]

The matrix $\mathbf{M}$ is
\[\mathbf{M} = \begin{pmatrix} a & 0.6 \\ 0.6 & b \end{pmatrix}\]
where $a$ and $b$ are constants with $a > 0$. \\

The matrix $\mathbf{M}$ represents a reflection in a line through the origin. \\
\begin{enumerate}[label=\textbf{(\alph*)}]
  \item Write down the eigenvalues of $\mathbf{M}$. \marks{1} \\
  \item By considering the characteristic equation of $\mathbf{M}$, use your answer to part (a) to find the values of $a$ and $b$. \marks{3} \\
  \item Hence find the equation of the mirror line of the reflection represented by $\mathbf{M}$. \marks{2} \\
\end{enumerate}

\vspace*{10pt}

\begin{djsSolution}
\begin{enumerate}[label=\textbf{(\alph*)}, start=1]
  \item A vector along the mirror line is unchanged by the reflection, and a vector perpendicular to the mirror line is reversed. So the eigenvalues are
\[
\lambda=1,\qquad \lambda=-1
\]

  \item The characteristic equation is \(\det(\mathbf{M}-\lambda\mathbf{I})=0\):
\begin{align*}
(a-\lambda)(b-\lambda)-0.6\times0.6&=0\\
\lambda^2-(a+b)\lambda+(ab-0.36)&=0
\end{align*}

By part (a), its roots are \(1\) and \(-1\), so it must be
\[
(\lambda-1)(\lambda+1)=\lambda^2-1=0
\]

Comparing coefficients:
\[
a+b=0,\qquad ab-0.36=-1
\]

So \(b=-a\) and \(ab=-0.64\), which gives
\[
-a^2=-0.64\quad\Longrightarrow\quad a^2=0.64
\]

Since \(a>0\), we obtain
\[
a=0.8,\qquad b=-0.8
\]

  \item The mirror line consists of the invariant points, so it is the line through the origin in the direction of an eigenvector with eigenvalue \(1\). Solving \((\mathbf{M}-\mathbf{I})\begin{pmatrix}x\\y\end{pmatrix}=\mathbf{0}\):
\[
\begin{pmatrix}-0.2 & 0.6\\ 0.6 & -1.8\end{pmatrix}\begin{pmatrix}x\\y\end{pmatrix}=\begin{pmatrix}0\\0\end{pmatrix}\quad\Longrightarrow\quad -0.2x+0.6y=0\quad\Longrightarrow\quad x=3y
\]

So an eigenvector is \(\begin{pmatrix}3\\1\end{pmatrix}\), and the mirror line is
\[
y=\dfrac{1}{3}x
\]

\end{enumerate}
\end{djsSolution}

\newpage

% ---- Question 3 ----
\questionitem
% [Source: _QBT___Solns__Eigenvalues__Eigenvectors_and_Diagonalisation, Q3]

The matrix $\mathbf{M}$ is given by
\[\mathbf{M} = \begin{pmatrix} 2 & k \\ k & 5 \end{pmatrix}\]
where $k$ is a negative constant. \\

It is given that $1$ is an eigenvalue of $\mathbf{M}$. \\

\begin{enumerate}[label=\textbf{(\alph*)}]
  \item Find the value of $k$. Hence find the other eigenvalue of $\mathbf{M}$, and a corresponding eigenvector for each eigenvalue of $\mathbf{M}$. \marks{6} \\
  \item
  \begin{enumerate}[label=(\roman*)]
    \item Describe how the directions of the invariant lines through the origin of the transformation represented by $\mathbf{M}$ are related to each other. Fully justify your answer. \marks{2}
    \item Describe fully the transformation represented by $\mathbf{M}$, stating the equation of any line of invariant points. \marks{3} \\
  \end{enumerate}
\end{enumerate}

\vspace*{10pt}

\begin{djsSolution}
\begin{enumerate}[label=\textbf{(\alph*)}, start=1]
  \item The eigenvalues satisfy \(\det(\mathbf{M}-\lambda\mathbf{I})=0\), which is
\[
\begin{vmatrix} 2-\lambda & k \\ k & 5-\lambda \end{vmatrix}=(2-\lambda)(5-\lambda)-k^2=0
\]

Since \(\lambda=1\) is an eigenvalue:
\[
1\times4-k^2=0\quad\Longrightarrow\quad k^2=4\quad\Longrightarrow\quad k=\pm 2
\]

Since \(k\) is negative, \(k=-2\).

The characteristic equation is then
\begin{align*}
(2-\lambda)(5-\lambda)-4&=0\\
\lambda^2-7\lambda+6&=0\\
(\lambda-1)(\lambda-6)&=0
\end{align*}

so the other eigenvalue is \(\lambda=6\).

For \(\lambda=1\), solve \((\mathbf{M}-\mathbf{I})\begin{pmatrix}x\\y\end{pmatrix}=\mathbf{0}\):
\[
\begin{pmatrix} 1 & -2 \\ -2 & 4 \end{pmatrix}\begin{pmatrix}x\\y\end{pmatrix}=\begin{pmatrix}0\\0\end{pmatrix}\quad\Longrightarrow\quad x-2y=0
\]

so a corresponding eigenvector is \(\begin{pmatrix}2\\1\end{pmatrix}\).

For \(\lambda=6\), solve \((\mathbf{M}-6\mathbf{I})\begin{pmatrix}x\\y\end{pmatrix}=\mathbf{0}\):
\[
\begin{pmatrix} -4 & -2 \\ -2 & -1 \end{pmatrix}\begin{pmatrix}x\\y\end{pmatrix}=\begin{pmatrix}0\\0\end{pmatrix}\quad\Longrightarrow\quad 2x+y=0
\]

so a corresponding eigenvector is \(\begin{pmatrix}1\\-2\end{pmatrix}\).

  \item
  \begin{enumerate}[label=(\roman*), start=1]
    \item The invariant lines through the origin are the lines in the directions of the eigenvectors, namely \(y=\frac{1}{2}x\) and \(y=-2x\). Their scalar product is
\[
\begin{pmatrix}2\\1\end{pmatrix}\cdot\begin{pmatrix}1\\-2\end{pmatrix}=2-2=0
\]

so the two invariant lines through the origin are perpendicular to each other.

    \item Using the eigenvectors from part (a):
\begin{align*}
\begin{pmatrix} 2 & -2 \\ -2 & 5 \end{pmatrix}\begin{pmatrix}2\\1\end{pmatrix}&=\begin{pmatrix}2\times2+(-2)\times1\\(-2)\times2+5\times1\end{pmatrix}=\begin{pmatrix}2\\1\end{pmatrix}\\
\begin{pmatrix} 2 & -2 \\ -2 & 5 \end{pmatrix}\begin{pmatrix}1\\-2\end{pmatrix}&=\begin{pmatrix}2\times1+(-2)\times(-2)\\(-2)\times1+5\times(-2)\end{pmatrix}=\begin{pmatrix}6\\-12\end{pmatrix}=6\begin{pmatrix}1\\-2\end{pmatrix}
\end{align*}

Any point can be written with position vector \(s\begin{pmatrix}2\\1\end{pmatrix}+t\begin{pmatrix}1\\-2\end{pmatrix}\), and \(\mathbf{M}\) maps it to
\[
s\begin{pmatrix}2\\1\end{pmatrix}+6t\begin{pmatrix}1\\-2\end{pmatrix}
\]

So every point on \(y=\frac{1}{2}x\) (where \(t=0\)) is fixed, while the component in the perpendicular direction \(\begin{pmatrix}1\\-2\end{pmatrix}\) is multiplied by \(6\).

Hence \(\mathbf{M}\) represents a one-way stretch with scale factor \(6\), parallel to the line \(y=-2x\), with line of invariant points \(y=\frac{1}{2}x\).
  \end{enumerate}
\end{enumerate}
\end{djsSolution}

\newpage

% ---- Question 4 ----
\questionitem
% [Source: _QBT___Solns__Eigenvalues__Eigenvectors_and_Diagonalisation, Q4]

A $3 \times 3$ matrix $\mathbf{A}$ has eigenvalues $2$, $-1$ and $1$, with corresponding eigenvectors
\[\begin{pmatrix} 1 \\ 0 \\ 1 \end{pmatrix}, \quad \begin{pmatrix} 0 \\ 1 \\ 1 \end{pmatrix}, \quad \begin{pmatrix} 1 \\ 1 \\ 1 \end{pmatrix}\]
respectively. \\
\begin{enumerate}[label=\textbf{(\alph*)}]
  \item By expressing $\mathbf{A}$ in the form $\mathbf{P}\mathbf{D}\mathbf{P}^{-1}$, where $\mathbf{D}$ is a diagonal matrix, find the matrix $\mathbf{A}$. \marks{4} \\
  \item Hence find the matrix $\mathbf{A}^4$. \marks{2} \\
\end{enumerate}

\vspace*{10pt}

\begin{djsSolution}
\begin{enumerate}[label=\textbf{(\alph*)}, start=1]
  \item Form \(\mathbf{P}\) with the eigenvectors as its columns, and \(\mathbf{D}\) with the corresponding eigenvalues on its leading diagonal in the same order:
\[
\mathbf{P}=\begin{pmatrix} 1 & 0 & 1 \\ 0 & 1 & 1 \\ 1 & 1 & 1 \end{pmatrix},\qquad \mathbf{D}=\begin{pmatrix} 2 & 0 & 0 \\ 0 & -1 & 0 \\ 0 & 0 & 1 \end{pmatrix}
\]

Then \(\mathbf{A}=\mathbf{P}\mathbf{D}\mathbf{P}^{-1}\).

Expanding along the first row, the determinant of \(\mathbf{P}\) is
\[
\det(\mathbf{P})=1\times(1-1)-0+1\times(0-1)=-1
\]

The inverse is given by
\[
\mathbf{P}^{-1}=\frac{1}{\det(\mathbf{P})}\operatorname{adj}(\mathbf{P})
\]

where \(\operatorname{adj}(\mathbf{P})\) is the transposed matrix of cofactors. The matrix of cofactors of \(\mathbf{P}\) is
\[
\begin{pmatrix} 0 & 1 & -1 \\ 1 & 0 & -1 \\ -1 & -1 & 1 \end{pmatrix}
\]

This matrix is symmetric, so it is equal to its transpose, and
\[
\mathbf{P}^{-1}=-\begin{pmatrix} 0 & 1 & -1 \\ 1 & 0 & -1 \\ -1 & -1 & 1 \end{pmatrix}=\begin{pmatrix} 0 & -1 & 1 \\ -1 & 0 & 1 \\ 1 & 1 & -1 \end{pmatrix}
\]

Multiplying \(\mathbf{P}\) by the diagonal matrix \(\mathbf{D}\) scales each column of \(\mathbf{P}\) by its eigenvalue:
\[
\mathbf{P}\mathbf{D}=\begin{pmatrix} 2 & 0 & 1 \\ 0 & -1 & 1 \\ 2 & -1 & 1 \end{pmatrix}
\]

Therefore
\begin{align*}
\mathbf{A}=\mathbf{P}\mathbf{D}\mathbf{P}^{-1}&=\begin{pmatrix} 2 & 0 & 1 \\ 0 & -1 & 1 \\ 2 & -1 & 1 \end{pmatrix}\begin{pmatrix} 0 & -1 & 1 \\ -1 & 0 & 1 \\ 1 & 1 & -1 \end{pmatrix}\\
&=\begin{pmatrix} 1 & -1 & 1 \\ 2 & 1 & -2 \\ 2 & -1 & 0 \end{pmatrix}
\end{align*}

  \item Using part (a):
\[
\mathbf{A}^4=\left(\mathbf{P}\mathbf{D}\mathbf{P}^{-1}\right)^4=\mathbf{P}\mathbf{D}\left(\mathbf{P}^{-1}\mathbf{P}\right)\mathbf{D}\left(\mathbf{P}^{-1}\mathbf{P}\right)\mathbf{D}\left(\mathbf{P}^{-1}\mathbf{P}\right)\mathbf{D}\mathbf{P}^{-1}=\mathbf{P}\mathbf{D}^4\mathbf{P}^{-1}
\]

where
\[
\mathbf{D}^4=\begin{pmatrix} 16 & 0 & 0 \\ 0 & 1 & 0 \\ 0 & 0 & 1 \end{pmatrix},\qquad \mathbf{P}\mathbf{D}^4=\begin{pmatrix} 16 & 0 & 1 \\ 0 & 1 & 1 \\ 16 & 1 & 1 \end{pmatrix}
\]

so
\begin{align*}
\mathbf{A}^4&=\begin{pmatrix} 16 & 0 & 1 \\ 0 & 1 & 1 \\ 16 & 1 & 1 \end{pmatrix}\begin{pmatrix} 0 & -1 & 1 \\ -1 & 0 & 1 \\ 1 & 1 & -1 \end{pmatrix}\\
&=\begin{pmatrix} 1 & -15 & 15 \\ 0 & 1 & 0 \\ 0 & -15 & 16 \end{pmatrix}
\end{align*}

\end{enumerate}
\end{djsSolution}

\newpage

% ---- Question 5 ----
\questionitem
% [Source: _QBT___Solns__Eigenvalues__Eigenvectors_and_Diagonalisation, Q5]

The matrix $\mathbf{M}$ is given by
\[\mathbf{M} = \begin{pmatrix} 4 & -1 & 1 \\ a & 5 & a \\ -8 & b & -5 \end{pmatrix}\]
where $a$ and $b$ are real numbers. \\

The vectors
\[\mathbf{v}_1 = \begin{pmatrix} 1 \\ 2 \\ 0 \end{pmatrix} \quad \text{and} \quad \mathbf{v}_2 = \begin{pmatrix} 0 \\ 1 \\ 1 \end{pmatrix}\]
are eigenvectors of $\mathbf{M}$, with corresponding eigenvalues $\lambda_1$ and $\lambda_2$ respectively. \\

\begin{enumerate}[label=\textbf{(\alph*)}]
  \item
  \begin{enumerate}[label=(\roman*)]
    \item Find the value of $\lambda_1$. \marks{1}
    \item Find the value of $a$ and the value of $b$. \marks{2} \\
  \end{enumerate}
\end{enumerate}
It is also given that $\lambda_3 = 3$ is an eigenvalue of $\mathbf{M}$. \\
\begin{enumerate}[label=\textbf{(\alph*)}, resume]
  \item Find the value of $\lambda_2$, and find an eigenvector $\mathbf{v}_3$ of $\mathbf{M}$ corresponding to the eigenvalue $\lambda_3 = 3$. \marks{4} \\
  \item Find matrices $\mathbf{U}$, $\mathbf{D}$ and $\mathbf{U}^{-1}$, such that $\mathbf{D}$ is diagonal and $\mathbf{M} = \mathbf{U}\mathbf{D}\mathbf{U}^{-1}$. \marks{3} \\
\end{enumerate}

\vspace*{10pt}

\begin{djsSolution}
\begin{enumerate}[label=\textbf{(\alph*)}, start=1]
  \item
  \begin{enumerate}[label=(\roman*), start=1]
    \item Since \(\mathbf{v}_1\) is an eigenvector, \(\mathbf{M}\mathbf{v}_1=\lambda_1\mathbf{v}_1\):
\[
\begin{pmatrix} 4 & -1 & 1 \\ a & 5 & a \\ -8 & b & -5 \end{pmatrix}\begin{pmatrix} 1 \\ 2 \\ 0 \end{pmatrix}=\begin{pmatrix} 2 \\ a+10 \\ 2b-8 \end{pmatrix}=\lambda_1\begin{pmatrix} 1 \\ 2 \\ 0 \end{pmatrix}
\]

Comparing the first components:
\[
\lambda_1=2
\]

    \item Comparing the second and third components, with \(\lambda_1=2\):
\begin{align*}
a+10&=4 \quad\Longrightarrow\quad a=-6\\
2b-8&=0 \quad\Longrightarrow\quad b=4
\end{align*}

  \end{enumerate}
  \item With \(a=-6\) and \(b=4\), the matrix is
\[
\mathbf{M}=\begin{pmatrix} 4 & -1 & 1 \\ -6 & 5 & -6 \\ -8 & 4 & -5 \end{pmatrix}
\]

Then
\[
\mathbf{M}\mathbf{v}_2=\begin{pmatrix} 4 & -1 & 1 \\ -6 & 5 & -6 \\ -8 & 4 & -5 \end{pmatrix}\begin{pmatrix} 0 \\ 1 \\ 1 \end{pmatrix}=\begin{pmatrix} 0 \\ -1 \\ -1 \end{pmatrix}=-1\begin{pmatrix} 0 \\ 1 \\ 1 \end{pmatrix}
\]

so \(\lambda_2=-1\).

For \(\lambda_3=3\), let \(\mathbf{v}_3=\begin{pmatrix} x \\ y \\ z \end{pmatrix}\) satisfy \((\mathbf{M}-3\mathbf{I})\mathbf{v}_3=\mathbf{0}\):
\[
\begin{pmatrix} 1 & -1 & 1 \\ -6 & 2 & -6 \\ -8 & 4 & -8 \end{pmatrix}\begin{pmatrix} x \\ y \\ z \end{pmatrix}=\begin{pmatrix} 0 \\ 0 \\ 0 \end{pmatrix}
\]

The first two rows give
\begin{align*}
x-y+z&=0\\
-6x+2y-6z&=0
\end{align*}

The first gives \(y=x+z\); substituting into the second:
\[
-6x+2(x+z)-6z=0 \quad\Longrightarrow\quad -4x-4z=0 \quad\Longrightarrow\quad z=-x
\]

and then \(y=0\). These also satisfy the third row, \(-8x+4y-8z=0\). Taking \(x=1\):
\[
\mathbf{v}_3=\begin{pmatrix} 1 \\ 0 \\ -1 \end{pmatrix}
\]

(or any non-zero multiple of this).

  \item Take the columns of \(\mathbf{U}\) to be the eigenvectors \(\mathbf{v}_1\), \(\mathbf{v}_2\), \(\mathbf{v}_3\), and the diagonal entries of \(\mathbf{D}\) to be the corresponding eigenvalues in the same order:
\[
\mathbf{U}=\begin{pmatrix} 1 & 0 & 1 \\ 2 & 1 & 0 \\ 0 & 1 & -1 \end{pmatrix},\qquad \mathbf{D}=\begin{pmatrix} 2 & 0 & 0 \\ 0 & -1 & 0 \\ 0 & 0 & 3 \end{pmatrix}
\]

The determinant is
\[
\det(\mathbf{U})=1\times(-1-0)-0+1\times(2-0)=1
\]

which is non-zero, so \(\mathbf{U}\) is invertible and
\[
\mathbf{U}^{-1}=\begin{pmatrix} -1 & 1 & -1 \\ 2 & -1 & 2 \\ 2 & -1 & 1 \end{pmatrix}
\]

Then \(\mathbf{M}=\mathbf{U}\mathbf{D}\mathbf{U}^{-1}\). (Any ordering of the eigenvectors in \(\mathbf{U}\), with the eigenvalues in \(\mathbf{D}\) in the matching order, is equally valid.)
\end{enumerate}
\end{djsSolution}

\newpage

% ---- Question 6 ----
\questionitem
% [Source: _QBT___Solns__Eigenvalues__Eigenvectors_and_Diagonalisation, Q6]

\[\mathbf{M} = \begin{pmatrix} 1 & -2 & 2 \\ -2 & 2 & 0 \\ 2 & 0 & 0 \end{pmatrix}\] \\
\begin{enumerate}[label=\textbf{(\alph*)}]
  \item Show that $-2$ is an eigenvalue of $\mathbf{M}$, and find a corresponding eigenvector. \marks{4} \\
  \item Given that $\mathbf{MU} = \mathbf{UD}$, where $\mathbf{D}$ is a diagonal matrix and each column of $\mathbf{U}$ is an eigenvector of $\mathbf{M}$ of unit length, find possible matrices for $\mathbf{D}$ and $\mathbf{U}$. \marks{11} \\
\end{enumerate}

\vspace*{10pt}

\begin{djsSolution}
\begin{enumerate}[label=\textbf{(\alph*)}, start=1]
  \item \(-2\) is an eigenvalue if \(\det(\mathbf{M}+2\mathbf{I})=0\). We have
\[
\mathbf{M}+2\mathbf{I}=\begin{pmatrix} 3 & -2 & 2 \\ -2 & 4 & 0 \\ 2 & 0 & 2 \end{pmatrix}
\]

Expanding along the top row:
\begin{align*}
\det(\mathbf{M}+2\mathbf{I}) &= 3\times(4\times 2-0\times 0)+2\times\big((-2)\times 2-0\times 2\big)+2\times\big((-2)\times 0-4\times 2\big)\\
&= 24-8-16\\
&= 0
\end{align*}

so \(-2\) is an eigenvalue of \(\mathbf{M}\).

For an eigenvector \(\begin{pmatrix}x\\y\\z\end{pmatrix}\), solve \((\mathbf{M}+2\mathbf{I})\begin{pmatrix}x\\y\\z\end{pmatrix}=\mathbf{0}\):
\begin{align*}
3x-2y+2z&=0\\
-2x+4y&=0\\
2x+2z&=0
\end{align*}

The second equation gives \(x=2y\) and the third gives \(z=-x\). Taking \(y=1\) gives \(x=2\), \(z=-2\), which also satisfies the first equation since \(6-2-4=0\). A corresponding eigenvector is
\[
\begin{pmatrix}2\\1\\-2\end{pmatrix}
\]

  \item The \(j\)th column of \(\mathbf{MU}\) is \(\mathbf{M}\mathbf{u}_j\), where \(\mathbf{u}_j\) is the \(j\)th column of \(\mathbf{U}\), and the \(j\)th column of \(\mathbf{UD}\) is \(d_j\mathbf{u}_j\), where \(d_j\) is the \(j\)th diagonal entry of \(\mathbf{D}\). So the diagonal entries of \(\mathbf{D}\) are the eigenvalues of \(\mathbf{M}\), in the same order as the corresponding eigenvectors in \(\mathbf{U}\).

The characteristic equation is \(\det(\mathbf{M}-\lambda\mathbf{I})=0\):
\begin{align*}
\begin{vmatrix} 1-\lambda & -2 & 2 \\ -2 & 2-\lambda & 0 \\ 2 & 0 & -\lambda \end{vmatrix}&=0\\
(1-\lambda)\big((2-\lambda)(-\lambda)-0\big)+2\big((-2)(-\lambda)-0\big)+2\big(0-2(2-\lambda)\big)&=0\\
(1-\lambda)(\lambda^2-2\lambda)+4\lambda-8+4\lambda&=0\\
-\lambda^3+3\lambda^2-2\lambda+8\lambda-8&=0\\
\lambda^3-3\lambda^2-6\lambda+8&=0
\end{align*}

From part (a), \(\lambda=-2\) is a root, so \((\lambda+2)\) is a factor:
\[
(\lambda+2)(\lambda^2-5\lambda+4)=0\quad\Longrightarrow\quad(\lambda+2)(\lambda-1)(\lambda-4)=0
\]

so the eigenvalues are \(\lambda=-2,\ 1,\ 4\).

\textbf{\(\lambda=1\):} solve \((\mathbf{M}-\mathbf{I})\begin{pmatrix}x\\y\\z\end{pmatrix}=\mathbf{0}\):
\begin{align*}
-2y+2z&=0\\
-2x+y&=0\\
2x-z&=0
\end{align*}

So \(y=z\) and \(y=2x\). Taking \(x=1\) gives the eigenvector \(\begin{pmatrix}1\\2\\2\end{pmatrix}\).

\textbf{\(\lambda=4\):} solve \((\mathbf{M}-4\mathbf{I})\begin{pmatrix}x\\y\\z\end{pmatrix}=\mathbf{0}\):
\begin{align*}
-3x-2y+2z&=0\\
-2x-2y&=0\\
2x-4z&=0
\end{align*}

So \(y=-x\) and \(x=2z\). Taking \(z=1\) gives \(x=2\), \(y=-2\), which satisfies the first equation since \(-6+4+2=0\). The eigenvector is \(\begin{pmatrix}2\\-2\\1\end{pmatrix}\).

Each of the three eigenvectors has length 3, for example
\[
\left|\begin{pmatrix}2\\1\\-2\end{pmatrix}\right|=\sqrt{4+1+4}=3
\]

and similarly \(\sqrt{1+4+4}=3\) and \(\sqrt{4+4+1}=3\). Dividing each by 3 gives unit eigenvectors, so possible matrices are
\[
\mathbf{D}=\begin{pmatrix} -2 & 0 & 0 \\ 0 & 1 & 0 \\ 0 & 0 & 4 \end{pmatrix},\qquad
\mathbf{U}=\dfrac{1}{3}\begin{pmatrix} 2 & 1 & 2 \\ 1 & 2 & -2 \\ -2 & 2 & 1 \end{pmatrix}
\]

The columns of \(\mathbf{U}\) may be listed in any order, provided the diagonal entries of \(\mathbf{D}\) are reordered to match, and any column may be multiplied by \(-1\).
\end{enumerate}
\end{djsSolution}

\newpage

% ---- Question 7 ----
\questionitem
% [Source: _QBT___Solns__Eigenvalues__Eigenvectors_and_Diagonalisation, Q7]

The matrix $\mathbf{A}$ is given by
\[\mathbf{A} = \begin{pmatrix} \dfrac{3}{8} & \dfrac{5}{16} \\[6pt] \dfrac{5}{4} & \dfrac{3}{8} \end{pmatrix}\] \\
\begin{enumerate}[label=\textbf{(\alph*)}]
  \item Find the eigenvalues of $\mathbf{A}$ and a corresponding eigenvector for each eigenvalue. \marks{5} \\
  \item Write down a matrix $\mathbf{P}$ and a diagonal matrix $\mathbf{D}$ such that $\mathbf{A} = \mathbf{P}\mathbf{D}\mathbf{P}^{-1}$. \marks{1} \\
  \item Find $\mathbf{A}^n$, where $n$ is a positive integer, as a single matrix whose entries are expressed in terms of $n$. \\
  Hence find the matrix $\mathbf{L}$ such that $\mathbf{A}^n \to \mathbf{L}$ as $n \to \infty$, explaining briefly why this limit exists. \marks{5} \\
  \item The transformation represented by $\mathbf{L}$ maps every point of the plane onto a single line through the origin. \\
  Find the equation of this line. \marks{2} \\
\end{enumerate}

\vspace*{10pt}

\begin{djsSolution}
\begin{enumerate}[label=\textbf{(\alph*)}, start=1]
  \item The eigenvalues satisfy \(\det(\mathbf{A}-\lambda\mathbf{I})=0\):
\begin{align*}
\left(\dfrac{3}{8}-\lambda\right)^2-\dfrac{5}{16}\times\dfrac{5}{4}&=0\\
\lambda^2-\dfrac{3}{4}\lambda+\dfrac{9}{64}-\dfrac{25}{64}&=0\\
4\lambda^2-3\lambda-1&=0\\
(4\lambda+1)(\lambda-1)&=0
\end{align*}

so the eigenvalues are \(\lambda=1\) and \(\lambda=-\frac{1}{4}\).

For \(\lambda=1\), \((\mathbf{A}-\mathbf{I})\begin{pmatrix}x\\y\end{pmatrix}=\mathbf{0}\) gives
\[
-\dfrac{5}{8}x+\dfrac{5}{16}y=0\quad\Longrightarrow\quad y=2x
\]

so an eigenvector is \(\begin{pmatrix}1\\2\end{pmatrix}\).

For \(\lambda=-\frac{1}{4}\), \(\left(\mathbf{A}+\frac14\mathbf{I}\right)\begin{pmatrix}x\\y\end{pmatrix}=\mathbf{0}\) gives
\[
\dfrac{5}{8}x+\dfrac{5}{16}y=0\quad\Longrightarrow\quad y=-2x
\]

so an eigenvector is \(\begin{pmatrix}1\\-2\end{pmatrix}\).

  \item Using the eigenvectors as columns, in the same order as the eigenvalues on the diagonal:
\[
\mathbf{P}=\begin{pmatrix}1&1\\2&-2\end{pmatrix},\qquad \mathbf{D}=\begin{pmatrix}1&0\\0&-\dfrac{1}{4}\end{pmatrix}
\]

  \item Since \(\mathbf{A}=\mathbf{P}\mathbf{D}\mathbf{P}^{-1}\), in the product \(\mathbf{A}^n=(\mathbf{P}\mathbf{D}\mathbf{P}^{-1})(\mathbf{P}\mathbf{D}\mathbf{P}^{-1})\cdots(\mathbf{P}\mathbf{D}\mathbf{P}^{-1})\) each inner \(\mathbf{P}^{-1}\mathbf{P}=\mathbf{I}\), so
\[
\mathbf{A}^n=\mathbf{P}\mathbf{D}^n\mathbf{P}^{-1},\qquad \mathbf{D}^n=\begin{pmatrix}1&0\\0&\left(-\dfrac{1}{4}\right)^n\end{pmatrix}
\]

Using
\[
\mathbf{P}^{-1}=\dfrac{1}{\det(\mathbf{P})}\operatorname{adj}(\mathbf{P})
\]

where \(\operatorname{adj}(\mathbf{P})\) is the transposed matrix of cofactors, with \(\det(\mathbf{P})=-2-2=-4\):
\[
\mathbf{P}^{-1}=-\dfrac{1}{4}\begin{pmatrix}-2&-1\\-2&1\end{pmatrix}=\begin{pmatrix}\dfrac{1}{2}&\dfrac{1}{4}\\[6pt]\dfrac{1}{2}&-\dfrac{1}{4}\end{pmatrix}
\]

Then
\[
\mathbf{P}\mathbf{D}^n=\begin{pmatrix}1&\left(-\dfrac{1}{4}\right)^n\\[6pt]2&-2\left(-\dfrac{1}{4}\right)^n\end{pmatrix}
\]

so
\[
\mathbf{A}^n=\begin{pmatrix}1&\left(-\dfrac{1}{4}\right)^n\\[6pt]2&-2\left(-\dfrac{1}{4}\right)^n\end{pmatrix}\begin{pmatrix}\dfrac{1}{2}&\dfrac{1}{4}\\[6pt]\dfrac{1}{2}&-\dfrac{1}{4}\end{pmatrix}=\begin{pmatrix}\dfrac{1}{2}+\dfrac{1}{2}\left(-\dfrac{1}{4}\right)^n&\dfrac{1}{4}-\dfrac{1}{4}\left(-\dfrac{1}{4}\right)^n\\[6pt]1-\left(-\dfrac{1}{4}\right)^n&\dfrac{1}{2}+\dfrac{1}{2}\left(-\dfrac{1}{4}\right)^n\end{pmatrix}
\]

Since \(\left|-\frac{1}{4}\right|<1\), \(\left(-\frac{1}{4}\right)^n\to0\) as \(n\to\infty\), so every entry of \(\mathbf{A}^n\) tends to a finite limit and the limit exists:
\[
\mathbf{L}=\begin{pmatrix}\dfrac{1}{2}&\dfrac{1}{4}\\[6pt]1&\dfrac{1}{2}\end{pmatrix}
\]

  \item The image of the point \((x,y)\) under \(\mathbf{L}\) is
\[
\begin{pmatrix}\dfrac{1}{2}&\dfrac{1}{4}\\[6pt]1&\dfrac{1}{2}\end{pmatrix}\begin{pmatrix}x\\y\end{pmatrix}=\begin{pmatrix}\dfrac{1}{2}x+\dfrac{1}{4}y\\[6pt]x+\dfrac{1}{2}y\end{pmatrix}=\left(\dfrac{1}{2}x+\dfrac{1}{4}y\right)\begin{pmatrix}1\\2\end{pmatrix}
\]

The second coordinate of every image point is twice the first, so the line is
\[
y=2x
\]

\end{enumerate}
\end{djsSolution}

\newpage

% ---- Question 8 ----
\questionitem
% [Source: _QBT___Solns__Eigenvalues__Eigenvectors_and_Diagonalisation, Q8]

The matrix $\mathbf{A}$ is given by
\[\mathbf{A} = \begin{pmatrix} -4 & -13 & 8 \\ -6 & -7 & 6 \\ -12 & -22 & 16 \end{pmatrix}\]
\begin{enumerate}[label=\textbf{(\alph*)}]
  \item Show that the characteristic equation of $\mathbf{A}$ is $\lambda^3 - 5\lambda^2 + 2\lambda + 8 = 0$ and hence find the eigenvalues of $\mathbf{A}$. \marks{4} \\
  \item Find a matrix $\mathbf{P}$ and a diagonal matrix $\mathbf{D}$ such that $\mathbf{A} = \mathbf{P}\mathbf{D}\mathbf{P}^{-1}$. \marks{5} \\
  \item Hence show that, for all positive integers $n$, the element in the first row and first column of $\mathbf{A}^n$ is $2(-1)^n - 2^n$. \marks{3} \\
\end{enumerate}

\vspace*{10pt}

\begin{djsSolution}
\begin{enumerate}[label=\textbf{(\alph*)}, start=1]
  \item The characteristic equation is \(\det(\mathbf{A}-\lambda\mathbf{I})=0\), where
\[
\det(\mathbf{A}-\lambda\mathbf{I})=\begin{vmatrix} -4-\lambda & -13 & 8 \\ -6 & -7-\lambda & 6 \\ -12 & -22 & 16-\lambda \end{vmatrix}
\]

Expanding along the first row:
\begin{align*}
\det(\mathbf{A}-\lambda\mathbf{I}) &= (-4-\lambda)\big[(-7-\lambda)(16-\lambda)+132\big]+13\big[-6(16-\lambda)+72\big]+8\big[132-12(7+\lambda)\big]\\
&= (-4-\lambda)(\lambda^2-9\lambda+20)+13(6\lambda-24)+8(48-12\lambda)\\
&= (-\lambda^3+5\lambda^2+16\lambda-80)+(78\lambda-312)+(384-96\lambda)\\
&= -\lambda^3+5\lambda^2-2\lambda-8
\end{align*}

Setting this equal to zero and multiplying by \(-1\), the characteristic equation is
\[
\lambda^3-5\lambda^2+2\lambda+8=0
\]

as required.

Substituting \(\lambda=-1\) gives \(-1-5-2+8=0\), so \((\lambda+1)\) is a factor:
\[
(\lambda+1)(\lambda^2-6\lambda+8)=(\lambda+1)(\lambda-2)(\lambda-4)=0
\]

The eigenvalues are therefore
\[
\lambda=-1,\qquad \lambda=2,\qquad \lambda=4
\]

  \item For each eigenvalue, solve \((\mathbf{A}-\lambda\mathbf{I})\mathbf{v}=\mathbf{0}\) with \(\mathbf{v}=\begin{pmatrix}x\\y\\z\end{pmatrix}\).

\(\lambda=-1\):
\[
\begin{pmatrix} -3 & -13 & 8 \\ -6 & -6 & 6 \\ -12 & -22 & 17 \end{pmatrix}\begin{pmatrix}x\\y\\z\end{pmatrix}=\mathbf{0}
\]

The second row gives \(-6x-6y+6z=0\), so \(z=x+y\). Substituting into the first row, \(-3x-13y+8(x+y)=5x-5y=0\), so \(y=x\) and \(z=2x\). Taking \(x=1\), an eigenvector is \(\begin{pmatrix}1\\1\\2\end{pmatrix}\).

\(\lambda=2\):
\[
\begin{pmatrix} -6 & -13 & 8 \\ -6 & -9 & 6 \\ -12 & -22 & 14 \end{pmatrix}\begin{pmatrix}x\\y\\z\end{pmatrix}=\mathbf{0}
\]

Subtracting the second equation from the first gives \(-4y+2z=0\), so \(z=2y\). The second row then gives \(-6x-9y+12y=0\), so \(y=2x\) and \(z=4x\). Taking \(x=1\), an eigenvector is \(\begin{pmatrix}1\\2\\4\end{pmatrix}\).

\(\lambda=4\):
\[
\begin{pmatrix} -8 & -13 & 8 \\ -6 & -11 & 6 \\ -12 & -22 & 12 \end{pmatrix}\begin{pmatrix}x\\y\\z\end{pmatrix}=\mathbf{0}
\]

Taking \(3\times\) the first equation minus \(4\times\) the second gives \(-39y+44y=5y=0\), so \(y=0\). The second row then gives \(-6x+6z=0\), so \(z=x\). Taking \(x=1\), an eigenvector is \(\begin{pmatrix}1\\0\\1\end{pmatrix}\).

Using the eigenvectors as the columns of \(\mathbf{P}\), in the same order as the eigenvalues on the diagonal of \(\mathbf{D}\):
\[
\mathbf{P}=\begin{pmatrix} 1 & 1 & 1 \\ 1 & 2 & 0 \\ 2 & 4 & 1 \end{pmatrix},\qquad \mathbf{D}=\begin{pmatrix} -1 & 0 & 0 \\ 0 & 2 & 0 \\ 0 & 0 & 4 \end{pmatrix}
\]

  \item Since \(\mathbf{P}^{-1}\mathbf{P}=\mathbf{I}\), the inner factors cancel in
\[
\mathbf{A}^n=(\mathbf{P}\mathbf{D}\mathbf{P}^{-1})(\mathbf{P}\mathbf{D}\mathbf{P}^{-1})\cdots(\mathbf{P}\mathbf{D}\mathbf{P}^{-1})=\mathbf{P}\mathbf{D}^n\mathbf{P}^{-1}
\]

where \(\mathbf{D}^n=\begin{pmatrix} (-1)^n & 0 & 0 \\ 0 & 2^n & 0 \\ 0 & 0 & 4^n \end{pmatrix}\). Here \(\det(\mathbf{P})=1\) and
\[
\mathbf{P}^{-1}=\begin{pmatrix} 2 & 3 & -2 \\ -1 & -1 & 1 \\ 0 & -2 & 1 \end{pmatrix}
\]

The first row of \(\mathbf{P}\mathbf{D}^n\) is \(\begin{pmatrix}(-1)^n & 2^n & 4^n\end{pmatrix}\), and the first column of \(\mathbf{P}^{-1}\) is \(\begin{pmatrix}2\\-1\\0\end{pmatrix}\). So the element in the first row and first column of \(\mathbf{A}^n\) is
\[
(-1)^n\times 2+2^n\times(-1)+4^n\times 0=2(-1)^n-2^n
\]

as required.
\end{enumerate}
\end{djsSolution}

\newpage

% ---- Question 9 ----
\questionitem
% [Source: _QBT___Solns__Eigenvalues__Eigenvectors_and_Diagonalisation, Q9]

The matrix $\mathbf{M}$ is given by
\[\mathbf{M} = \begin{pmatrix} 3 & 4 & -4 \\ 1 & 3 & -1 \\ 1 & 4 & -2 \end{pmatrix}\]
\begin{enumerate}[label=\textbf{(\alph*)}]
  \item Show that the characteristic equation of $\mathbf{M}$ is
  \[\lambda^3 - 4\lambda^2 + \lambda + 6 = 0\] \marks[-20pt]{3}
  \item Hence find the three eigenvalues of $\mathbf{M}$. \marks{2} \\
  \item For each eigenvalue of $\mathbf{M}$ find an associated eigenvector. \marks{4} \\
  \item Write down a matrix $\mathbf{P}$ and a diagonal matrix $\mathbf{D}$ such that $\mathbf{M} = \mathbf{P}\mathbf{D}\mathbf{P}^{-1}$. Hence find $\mathbf{M}^n$ as a single matrix, where $n$ is a positive integer, giving each entry in terms of $n$. \marks{4} \\
\end{enumerate}

\vspace*{10pt}

\begin{djsSolution}
\begin{enumerate}[label=\textbf{(\alph*)}, start=1]
  \item The characteristic equation is \(\det(\mathbf{M}-\lambda\mathbf{I})=0\), where
\[
\det(\mathbf{M}-\lambda\mathbf{I})=\begin{vmatrix} 3-\lambda & 4 & -4 \\ 1 & 3-\lambda & -1 \\ 1 & 4 & -2-\lambda \end{vmatrix}
\]

Expanding along the first row:
\begin{align*}
\det(\mathbf{M}-\lambda\mathbf{I}) &= (3-\lambda)\big[(3-\lambda)(-2-\lambda)+4\big]-4\big[(-2-\lambda)+1\big]-4\big[4-(3-\lambda)\big]\\
&= (3-\lambda)(\lambda^2-\lambda-2)-4(-1-\lambda)-4(1+\lambda)\\
&= (3-\lambda)(\lambda^2-\lambda-2)+4+4\lambda-4-4\lambda\\
&= 3\lambda^2-3\lambda-6-\lambda^3+\lambda^2+2\lambda\\
&= -\lambda^3+4\lambda^2-\lambda-6
\end{align*}

Setting this equal to zero and multiplying by \(-1\) gives
\[
\lambda^3-4\lambda^2+\lambda+6=0
\]

as required.

  \item Let \(f(\lambda)=\lambda^3-4\lambda^2+\lambda+6\). Then
\[
f(-1)=-1-4-1+6=0
\]

so, by the factor theorem, \((\lambda+1)\) is a factor. Factorising:
\[
\lambda^3-4\lambda^2+\lambda+6=(\lambda+1)(\lambda^2-5\lambda+6)=(\lambda+1)(\lambda-2)(\lambda-3)
\]

Hence the eigenvalues are \(\lambda=-1\), \(\lambda=2\) and \(\lambda=3\).

  \item Let \(\mathbf{v}=\begin{pmatrix}x\\y\\z\end{pmatrix}\) and solve \((\mathbf{M}-\lambda\mathbf{I})\mathbf{v}=\mathbf{0}\) for each eigenvalue.

\(\lambda=-1\):
\[
\begin{pmatrix} 4 & 4 & -4 \\ 1 & 4 & -1 \\ 1 & 4 & -1 \end{pmatrix}\begin{pmatrix}x\\y\\z\end{pmatrix}=\mathbf{0}
\quad\Longrightarrow\quad x+y-z=0,\qquad x+4y-z=0
\]

Subtracting gives \(3y=0\), so \(y=0\) and \(x=z\). An eigenvector is \(\begin{pmatrix}1\\0\\1\end{pmatrix}\).

\(\lambda=2\):
\[
\begin{pmatrix} 1 & 4 & -4 \\ 1 & 1 & -1 \\ 1 & 4 & -4 \end{pmatrix}\begin{pmatrix}x\\y\\z\end{pmatrix}=\mathbf{0}
\quad\Longrightarrow\quad x+4y-4z=0,\qquad x+y-z=0
\]

Subtracting gives \(3y-3z=0\), so \(y=z\) and then \(x=0\). An eigenvector is \(\begin{pmatrix}0\\1\\1\end{pmatrix}\).

\(\lambda=3\):
\[
\begin{pmatrix} 0 & 4 & -4 \\ 1 & 0 & -1 \\ 1 & 4 & -5 \end{pmatrix}\begin{pmatrix}x\\y\\z\end{pmatrix}=\mathbf{0}
\quad\Longrightarrow\quad 4y-4z=0,\qquad x-z=0,\qquad x+4y-5z=0
\]

So \(y=z\) and \(x=z\), which also satisfy the third equation. An eigenvector is \(\begin{pmatrix}1\\1\\1\end{pmatrix}\).

Any non-zero multiple of each of these vectors is also an eigenvector.

  \item Taking the eigenvectors as the columns of \(\mathbf{P}\), with the eigenvalues in the same order on the diagonal of \(\mathbf{D}\):
\[
\mathbf{P}=\begin{pmatrix} 1 & 0 & 1 \\ 0 & 1 & 1 \\ 1 & 1 & 1 \end{pmatrix},\qquad
\mathbf{D}=\begin{pmatrix} -1 & 0 & 0 \\ 0 & 2 & 0 \\ 0 & 0 & 3 \end{pmatrix}
\]

Then \(\det(\mathbf{P})=1\times(1-1)-0+1\times(0-1)=-1\neq 0\), and
\[
\mathbf{P}^{-1}=\begin{pmatrix} 0 & -1 & 1 \\ -1 & 0 & 1 \\ 1 & 1 & -1 \end{pmatrix}
\]

Since \(\mathbf{P}^{-1}\mathbf{P}=\mathbf{I}\), we have
\[
\mathbf{M}^n=(\mathbf{P}\mathbf{D}\mathbf{P}^{-1})(\mathbf{P}\mathbf{D}\mathbf{P}^{-1})\cdots(\mathbf{P}\mathbf{D}\mathbf{P}^{-1})=\mathbf{P}\mathbf{D}^n\mathbf{P}^{-1},\qquad
\mathbf{D}^n=\begin{pmatrix} (-1)^n & 0 & 0 \\ 0 & 2^n & 0 \\ 0 & 0 & 3^n \end{pmatrix}
\]

First calculate:
\[
\mathbf{P}\mathbf{D}^n=\begin{pmatrix} 1 & 0 & 1 \\ 0 & 1 & 1 \\ 1 & 1 & 1 \end{pmatrix}\begin{pmatrix} (-1)^n & 0 & 0 \\ 0 & 2^n & 0 \\ 0 & 0 & 3^n \end{pmatrix}=\begin{pmatrix} (-1)^n & 0 & 3^n \\ 0 & 2^n & 3^n \\ (-1)^n & 2^n & 3^n \end{pmatrix}
\]

Then
\begin{align*}
\mathbf{M}^n&=\begin{pmatrix} (-1)^n & 0 & 3^n \\ 0 & 2^n & 3^n \\ (-1)^n & 2^n & 3^n \end{pmatrix}\begin{pmatrix} 0 & -1 & 1 \\ -1 & 0 & 1 \\ 1 & 1 & -1 \end{pmatrix}\\
&=\begin{pmatrix} 3^n & 3^n-(-1)^n & (-1)^n-3^n \\ 3^n-2^n & 3^n & 2^n-3^n \\ 3^n-2^n & 3^n-(-1)^n & (-1)^n+2^n-3^n \end{pmatrix}
\end{align*}

\end{enumerate}
\end{djsSolution}

\newpage

% ---- Question 10 ----
\questionitem
% [Source: _QBT___Solns__Eigenvalues__Eigenvectors_and_Diagonalisation, Q10]

The square matrix $\mathbf{A}$ has an eigenvalue $\lambda$ with corresponding eigenvector $\mathbf{e}$. The non-singular matrix $\mathbf{M}$ is of the same order as $\mathbf{A}$, and the matrix $\mathbf{B}$ is defined by $\mathbf{B} = \mathbf{M}^{-1}\mathbf{A}\mathbf{M}$. \\
\begin{enumerate}[label=\textbf{(\alph*)}]
  \item Show that $\mathbf{M}^{-1}\mathbf{e}$ is an eigenvector of $\mathbf{B}$, and that $\lambda$ is the corresponding eigenvalue. \marks{3} \\
\end{enumerate}
Let
\[\mathbf{A} = \begin{pmatrix} 3 & 0 & 0 \\ 1 & -2 & 0 \\ 2 & 5 & 1 \end{pmatrix}\]
\begin{enumerate}[label=\textbf{(\alph*)}, resume]
  \item Write down the eigenvalues of $\mathbf{A}$ and obtain corresponding eigenvectors. \marks{4} \\
\end{enumerate}
It is given that
\[\mathbf{M} = \begin{pmatrix} 1 & 1 & 0 \\ 0 & 1 & 0 \\ 0 & 2 & 1 \end{pmatrix}\]
\begin{enumerate}[label=\textbf{(\alph*)}, resume]
  \item Find $\mathbf{M}^{-1}$ and hence, without finding $\mathbf{B}$ explicitly, find the eigenvalues and corresponding eigenvectors of $\mathbf{B}$. \marks{5} \\
\end{enumerate}

\vspace*{10pt}

\begin{djsSolution}
\begin{enumerate}[label=\textbf{(\alph*)}, start=1]
  \item Since \(\mathbf{e}\) is an eigenvector of \(\mathbf{A}\) with eigenvalue \(\lambda\), \(\mathbf{A}\mathbf{e}=\lambda\mathbf{e}\) with \(\mathbf{e}\neq\mathbf{0}\). Then
\begin{align*}
\mathbf{B}\left(\mathbf{M}^{-1}\mathbf{e}\right) &= \mathbf{M}^{-1}\mathbf{A}\mathbf{M}\mathbf{M}^{-1}\mathbf{e}\\
&= \mathbf{M}^{-1}\mathbf{A}\mathbf{I}\mathbf{e}\\
&= \mathbf{M}^{-1}\left(\lambda\mathbf{e}\right)\\
&= \lambda\left(\mathbf{M}^{-1}\mathbf{e}\right)
\end{align*}

Also \(\mathbf{M}^{-1}\mathbf{e}\neq\mathbf{0}\): if \(\mathbf{M}^{-1}\mathbf{e}=\mathbf{0}\), then multiplying on the left by \(\mathbf{M}\) would give \(\mathbf{e}=\mathbf{M}\mathbf{0}=\mathbf{0}\), which is false. Hence \(\mathbf{M}^{-1}\mathbf{e}\) is an eigenvector of \(\mathbf{B}\) with corresponding eigenvalue \(\lambda\), as required.

  \item \(\mathbf{A}\) is lower triangular, so
\[
\det(\mathbf{A}-\lambda\mathbf{I})=(3-\lambda)(-2-\lambda)(1-\lambda)
\]

and the eigenvalues are the diagonal entries, \(\lambda=3,\ -2,\ 1\).

For an eigenvector \(\begin{pmatrix}x\\y\\z\end{pmatrix}\), solve \((\mathbf{A}-\lambda\mathbf{I})\begin{pmatrix}x\\y\\z\end{pmatrix}=\mathbf{0}\).

\(\lambda=3\):
\[
\begin{pmatrix}0&0&0\\1&-5&0\\2&5&-2\end{pmatrix}\begin{pmatrix}x\\y\\z\end{pmatrix}=\mathbf{0}\quad\Longrightarrow\quad x=5y,\qquad 2x+5y-2z=0
\]

Taking \(y=2\) gives \(x=10\) and \(2z=20+10=30\), so \(z=15\). An eigenvector is \(\begin{pmatrix}10\\2\\15\end{pmatrix}\).

\(\lambda=-2\):
\[
\begin{pmatrix}5&0&0\\1&0&0\\2&5&3\end{pmatrix}\begin{pmatrix}x\\y\\z\end{pmatrix}=\mathbf{0}\quad\Longrightarrow\quad x=0,\qquad 5y+3z=0
\]

Taking \(y=3\) gives \(z=-5\). An eigenvector is \(\begin{pmatrix}0\\3\\-5\end{pmatrix}\).

\(\lambda=1\):
\[
\begin{pmatrix}2&0&0\\1&-3&0\\2&5&0\end{pmatrix}\begin{pmatrix}x\\y\\z\end{pmatrix}=\mathbf{0}\quad\Longrightarrow\quad x=0,\qquad y=0
\]

with \(z\) free. An eigenvector is \(\begin{pmatrix}0\\0\\1\end{pmatrix}\).

Any non-zero multiples of these vectors are also eigenvectors.

  \item Expanding along the first column:
\[
\det(\mathbf{M})=1\times(1\times1-0\times2)=1
\]

Using
\[
\mathbf{M}^{-1}=\frac{1}{\det(\mathbf{M})}\operatorname{adj}(\mathbf{M})
\]

where \(\operatorname{adj}(\mathbf{M})\) is the transposed matrix of cofactors. The matrix of cofactors is
\[
\begin{pmatrix}1&0&0\\-1&1&-2\\0&0&1\end{pmatrix}
\]

Transposing:
\[
\mathbf{M}^{-1}=\begin{pmatrix}1&-1&0\\0&1&0\\0&-2&1\end{pmatrix}
\]

By part (a), each eigenvalue \(\lambda\) of \(\mathbf{A}\) with eigenvector \(\mathbf{e}\) is an eigenvalue of \(\mathbf{B}\) with eigenvector \(\mathbf{M}^{-1}\mathbf{e}\). Since \(\mathbf{B}\) is \(3\times3\) it has at most three eigenvalues, so the eigenvalues of \(\mathbf{B}\) are \(3,\ -2,\ 1\).

\(\lambda=3\):
\[
\begin{pmatrix}1&-1&0\\0&1&0\\0&-2&1\end{pmatrix}\begin{pmatrix}10\\2\\15\end{pmatrix}=\begin{pmatrix}10-2\\2\\-4+15\end{pmatrix}=\begin{pmatrix}8\\2\\11\end{pmatrix}
\]

\(\lambda=-2\):
\[
\begin{pmatrix}1&-1&0\\0&1&0\\0&-2&1\end{pmatrix}\begin{pmatrix}0\\3\\-5\end{pmatrix}=\begin{pmatrix}-3\\3\\-6-5\end{pmatrix}=\begin{pmatrix}-3\\3\\-11\end{pmatrix}
\]

\(\lambda=1\):
\[
\begin{pmatrix}1&-1&0\\0&1&0\\0&-2&1\end{pmatrix}\begin{pmatrix}0\\0\\1\end{pmatrix}=\begin{pmatrix}0\\0\\1\end{pmatrix}
\]

So \(\mathbf{B}\) has eigenvalue \(3\) with eigenvector \(\begin{pmatrix}8\\2\\11\end{pmatrix}\), eigenvalue \(-2\) with eigenvector \(\begin{pmatrix}-3\\3\\-11\end{pmatrix}\), and eigenvalue \(1\) with eigenvector \(\begin{pmatrix}0\\0\\1\end{pmatrix}\), or any non-zero multiples of these.
\end{enumerate}
\end{djsSolution}

\newpage

% ---- Question 11 ----
\questionitem
% [Source: _QBT___Solns__Eigenvalues__Eigenvectors_and_Diagonalisation, Q11]

It is given that $\lambda$ is an eigenvalue of the non-singular square matrix $\mathbf{A}$, with corresponding eigenvector $\mathbf{e}$. \\
\begin{enumerate}[label=\textbf{(\alph*)}]
  \item Show that $\lambda^{2}$ is an eigenvalue of $\mathbf{A}^{2}$ and that $\lambda^{-1}$ is an eigenvalue of $\mathbf{A}^{-1}$, and that $\mathbf{e}$ is a corresponding eigenvector in each case. \marks{3} \\
  \item Deduce that $\lambda^{2} + 2\lambda^{-1}$ is an eigenvalue of $\mathbf{A}^{2} + 2\mathbf{A}^{-1}$. \marks{1} \\
\end{enumerate}
It is given that $-1$ is an eigenvalue of the matrix $\mathbf{A}$, where
\[\mathbf{A} = \begin{pmatrix} 2 & 1 & -1 \\ 3 & 2 & -3 \\ 3 & 1 & -2 \end{pmatrix}\]
\begin{enumerate}[label=\textbf{(\alph*)}, resume]
  \item Find a corresponding eigenvector. \marks{2} \\
  \item It is also given that
  \[\begin{pmatrix} 1 \\ 0 \\ 1 \end{pmatrix} \quad \text{and} \quad \begin{pmatrix} 1 \\ 1 \\ 1 \end{pmatrix}\]
  are eigenvectors of the matrix $\mathbf{A}$. Find the corresponding eigenvalues. \marks{2} \\
  \item Hence find a matrix $\mathbf{P}$ and a diagonal matrix $\mathbf{D}$ such that
  \[\left(\mathbf{A}^{2} + 2\mathbf{A}^{-1}\right)^2 = \mathbf{P}\mathbf{D}\mathbf{P}^{-1}\] \marks[-20pt]{4}
\end{enumerate}

\vspace*{10pt}

\begin{djsSolution}
\begin{enumerate}[label=\textbf{(\alph*)}, start=1]
  \item Since \(\mathbf{e}\) is an eigenvector of \(\mathbf{A}\) with eigenvalue \(\lambda\), we have \(\mathbf{A}\mathbf{e}=\lambda\mathbf{e}\) with \(\mathbf{e}\neq\mathbf{0}\). Then
\[
\mathbf{A}^2\mathbf{e}=\mathbf{A}(\mathbf{A}\mathbf{e})=\mathbf{A}(\lambda\mathbf{e})=\lambda\mathbf{A}\mathbf{e}=\lambda^2\mathbf{e}
\]

so \(\lambda^2\) is an eigenvalue of \(\mathbf{A}^2\) with corresponding eigenvector \(\mathbf{e}\).

If \(\lambda=0\), then \(\mathbf{A}\mathbf{e}=\mathbf{0}\), so \(\mathbf{e}=\mathbf{A}^{-1}\mathbf{0}=\mathbf{0}\), a contradiction. Hence \(\lambda\neq0\). Premultiplying \(\mathbf{A}\mathbf{e}=\lambda\mathbf{e}\) by \(\mathbf{A}^{-1}\):
\[
\mathbf{e}=\lambda\mathbf{A}^{-1}\mathbf{e}\quad\Longrightarrow\quad \mathbf{A}^{-1}\mathbf{e}=\lambda^{-1}\mathbf{e}
\]

so \(\lambda^{-1}\) is an eigenvalue of \(\mathbf{A}^{-1}\) with corresponding eigenvector \(\mathbf{e}\).

  \item Using part (a):
\[
\left(\mathbf{A}^2+2\mathbf{A}^{-1}\right)\mathbf{e}=\mathbf{A}^2\mathbf{e}+2\mathbf{A}^{-1}\mathbf{e}=\lambda^2\mathbf{e}+2\lambda^{-1}\mathbf{e}=\left(\lambda^2+2\lambda^{-1}\right)\mathbf{e}
\]

so \(\lambda^2+2\lambda^{-1}\) is an eigenvalue of \(\mathbf{A}^2+2\mathbf{A}^{-1}\), as required.

  \item Let the eigenvector be \(\begin{pmatrix}x\\y\\z\end{pmatrix}\). It satisfies \((\mathbf{A}+\mathbf{I})\begin{pmatrix}x\\y\\z\end{pmatrix}=\mathbf{0}\), that is
\[
\begin{pmatrix} 3 & 1 & -1 \\ 3 & 3 & -3 \\ 3 & 1 & -1 \end{pmatrix}\begin{pmatrix}x\\y\\z\end{pmatrix}=\begin{pmatrix}0\\0\\0\end{pmatrix}
\]

This gives \(3x+y-z=0\) and \(3x+3y-3z=0\), i.e.\ \(x+y-z=0\). Subtracting gives \(2x=0\), so \(x=0\) and \(y=z\). A corresponding eigenvector is
\[
\begin{pmatrix}0\\1\\1\end{pmatrix}
\]

  \item Multiplying each eigenvector by \(\mathbf{A}\):
\begin{align*}
\begin{pmatrix} 2 & 1 & -1 \\ 3 & 2 & -3 \\ 3 & 1 & -2 \end{pmatrix}\begin{pmatrix}1\\0\\1\end{pmatrix}&=\begin{pmatrix}1\\0\\1\end{pmatrix}=1\begin{pmatrix}1\\0\\1\end{pmatrix}\\
\begin{pmatrix} 2 & 1 & -1 \\ 3 & 2 & -3 \\ 3 & 1 & -2 \end{pmatrix}\begin{pmatrix}1\\1\\1\end{pmatrix}&=\begin{pmatrix}2\\2\\2\end{pmatrix}=2\begin{pmatrix}1\\1\\1\end{pmatrix}
\end{align*}

The corresponding eigenvalues are \(1\) and \(2\) respectively.

  \item Let \(\mathbf{B}=\mathbf{A}^2+2\mathbf{A}^{-1}\). By part (b), each eigenvector of \(\mathbf{A}\) with eigenvalue \(\lambda\) is an eigenvector of \(\mathbf{B}\) with eigenvalue \(\mu=\lambda^2+2\lambda^{-1}\):
\begin{align*}
\lambda=-1:&\quad \mu=1-2=-1\\
\lambda=1:&\quad \mu=1+2=3\\
\lambda=2:&\quad \mu=4+1=5
\end{align*}

If \(\mathbf{B}\mathbf{e}=\mu\mathbf{e}\), then, exactly as in part (a), \(\mathbf{B}^2\mathbf{e}=\mathbf{B}(\mu\mathbf{e})=\mu\mathbf{B}\mathbf{e}=\mu^2\mathbf{e}\). So \(\mathbf{B}^2\) has the same eigenvectors with eigenvalues \((-1)^2=1\), \(3^2=9\) and \(5^2=25\).

Taking the eigenvectors as the columns of \(\mathbf{P}\) and the corresponding eigenvalues, in the same order, on the diagonal of \(\mathbf{D}\):
\[
\mathbf{P}=\begin{pmatrix} 0 & 1 & 1 \\ 1 & 0 & 1 \\ 1 & 1 & 1 \end{pmatrix},\qquad \mathbf{D}=\begin{pmatrix} 1 & 0 & 0 \\ 0 & 9 & 0 \\ 0 & 0 & 25 \end{pmatrix}
\]

Since \(\det(\mathbf{P})=0\times(0-1)-1\times(1-1)+1\times(1-0)=1\neq0\), \(\mathbf{P}\) is non-singular, so \(\left(\mathbf{A}^2+2\mathbf{A}^{-1}\right)^2=\mathbf{P}\mathbf{D}\mathbf{P}^{-1}\) with these matrices.
\end{enumerate}
\end{djsSolution}

\newpage

% ---- Question 12 ----
\questionitem
% [Source: _QBT___Solns__Eigenvalues__Eigenvectors_and_Diagonalisation, Q12]

The matrix $\mathbf{M}$ is given by
\[\mathbf{M} = \dfrac{1}{3}\begin{pmatrix} 2 & -1 & 2 \\ 2 & 2 & -1 \\ -1 & 2 & 2 \end{pmatrix}\] \\
\begin{enumerate}[label=\textbf{(\alph*)}]
  \item
  \begin{enumerate}[label=(\roman*)]
    \item Find the characteristic equation of $\mathbf{M}$, giving your answer in the form $\lambda^3 + p\lambda^2 + q\lambda + r = 0$, where $p$, $q$ and $r$ are integers. Hence show that $1$ is an eigenvalue of $\mathbf{M}$. \marks{3}
    \item Find the other two eigenvalues of $\mathbf{M}$, giving each in the form $x + yi$ where $x$ and $y$ are real. Deduce that $1$ is the only real eigenvalue of $\mathbf{M}$. \marks{3} \\
  \end{enumerate}
  \item A student claims that $\mathbf{M}$ can be written in the form $\mathbf{M} = \mathbf{P}\mathbf{D}\mathbf{P}^{-1}$, where $\mathbf{D}$ is a diagonal matrix and \textbf{both} $\mathbf{P}$ and $\mathbf{D}$ have only real entries. Using your answer to part (a)(ii), explain why the student is wrong. \marks{2} \\
\end{enumerate}
You are now given the following information.
\begin{itemize}
  \item $\mathbf{M}$ represents a rotation in 3-D space by an acute angle about an axis through the origin.
  \item $\mathbf{e}$ is a real eigenvector of $\mathbf{M}$ associated with the eigenvalue $1$.
\end{itemize}
\begin{enumerate}[label=\textbf{(\alph*)}, resume]
  \item
  \begin{enumerate}[label=(\roman*)]
    \item Find a possible vector $\mathbf{e}$. \marks{2}
    \item Explain the significance of $\mathbf{e}$, and the fact that its associated eigenvalue is $1$, in relation to the rotation that $\mathbf{M}$ represents. Hence write down a vector equation for the axis of this rotation. \marks{2}
    \item Explain \textbf{geometrically} why $\mathbf{M}$ has no real eigenvectors other than non-zero multiples of $\mathbf{e}$. \marks{1} \\
  \end{enumerate}
\end{enumerate}

\vspace*{10pt}

\begin{djsSolution}
\begin{enumerate}[label=\textbf{(\alph*)}, start=1]
  \item
  \begin{enumerate}[label=(\roman*), start=1]
    \item Since \(\mathbf{M}-\lambda\mathbf{I}=\frac{1}{3}\left(\mathbf{A}-3\lambda\mathbf{I}\right)\), where \(\mathbf{A}=\begin{pmatrix} 2 & -1 & 2 \\ 2 & 2 & -1 \\ -1 & 2 & 2 \end{pmatrix}\), and the matrices are \(3\times3\), we have
\[
\det(\mathbf{M}-\lambda\mathbf{I})=\dfrac{1}{27}\det\begin{pmatrix} 2-3\lambda & -1 & 2 \\ 2 & 2-3\lambda & -1 \\ -1 & 2 & 2-3\lambda \end{pmatrix}
\]

Expanding the determinant along the first row:
\begin{align*}
&(2-3\lambda)\left[(2-3\lambda)^2+2\right]+\left[2(2-3\lambda)-1\right]+2\left[4+(2-3\lambda)\right]\\
&=(2-3\lambda)^3+6(2-3\lambda)+7\\
&=8-36\lambda+54\lambda^2-27\lambda^3+12-18\lambda+7\\
&=-27\lambda^3+54\lambda^2-54\lambda+27
\end{align*}

So \(\det(\mathbf{M}-\lambda\mathbf{I})=-\lambda^3+2\lambda^2-2\lambda+1\), and the characteristic equation \(\det(\mathbf{M}-\lambda\mathbf{I})=0\) is
\[
\lambda^3-2\lambda^2+2\lambda-1=0
\]

with \(p=-2\), \(q=2\), \(r=-1\).

Substituting \(\lambda=1\):
\[
1-2+2-1=0
\]

so \(\lambda=1\) satisfies the characteristic equation and \(1\) is an eigenvalue of \(\mathbf{M}\).

    \item Since \(\lambda-1\) is a factor:
\[
\lambda^3-2\lambda^2+2\lambda-1=(\lambda-1)(\lambda^2-\lambda+1)
\]

The other eigenvalues satisfy \(\lambda^2-\lambda+1=0\), so
\[
\lambda=\dfrac{1\pm\sqrt{1-4}}{2}=\dfrac{1}{2}\pm\dfrac{\sqrt{3}}{2}i
\]

The discriminant of \(\lambda^2-\lambda+1\) is \(b^2-4ac=1-4=-3<0\), so this quadratic has no real roots. Hence the only real root of the characteristic equation, and so the only real eigenvalue of \(\mathbf{M}\), is \(1\).
  \end{enumerate}

  \item If \(\mathbf{M}=\mathbf{P}\mathbf{D}\mathbf{P}^{-1}\), then
\[
\det(\mathbf{M}-\lambda\mathbf{I})=\det\left(\mathbf{P}(\mathbf{D}-\lambda\mathbf{I})\mathbf{P}^{-1}\right)=\det(\mathbf{P})\det(\mathbf{D}-\lambda\mathbf{I})\det(\mathbf{P}^{-1})=\det(\mathbf{D}-\lambda\mathbf{I})
\]

so \(\mathbf{M}\) and \(\mathbf{D}\) have the same eigenvalues. The eigenvalues of the diagonal matrix \(\mathbf{D}\) are its diagonal entries, so these would have to be \(1\), \(\frac{1}{2}+\frac{\sqrt{3}}{2}i\) and \(\frac{1}{2}-\frac{\sqrt{3}}{2}i\). By part (a)(ii), two of these are not real, so \(\mathbf{D}\) cannot have only real entries, and the student is wrong.

  \item
  \begin{enumerate}[label=(\roman*), start=1]
    \item Let \(\mathbf{e}=\begin{pmatrix}x\\y\\z\end{pmatrix}\) satisfy \((\mathbf{M}-\mathbf{I})\mathbf{e}=\mathbf{0}\), that is
\[
\dfrac{1}{3}\begin{pmatrix} -1 & -1 & 2 \\ 2 & -1 & -1 \\ -1 & 2 & -1 \end{pmatrix}\begin{pmatrix}x\\y\\z\end{pmatrix}=\begin{pmatrix}0\\0\\0\end{pmatrix}
\]

The first two rows give \(-x-y+2z=0\) and \(2x-y-z=0\). From the second, \(y=2x-z\); substituting into the first:
\[
-x-(2x-z)+2z=0\quad\Longrightarrow\quad x=z\quad\Longrightarrow\quad y=x
\]

These also satisfy the third row, \(-x+2y-z=0\). Taking \(x=1\):
\[
\mathbf{e}=\begin{pmatrix}1\\1\\1\end{pmatrix}
\]

    \item Since \(\mathbf{M}\mathbf{e}=1\,\mathbf{e}=\mathbf{e}\), the vector \(\mathbf{e}\), and every multiple \(t\mathbf{e}\), is mapped to itself: it is neither turned nor stretched nor reversed. So every point on the line through the origin in the direction of \(\mathbf{e}\) is invariant under the rotation. The points left fixed by a rotation are exactly those on its axis, so this line is the axis of rotation:
\[
\mathbf{r}=t\begin{pmatrix}1\\1\\1\end{pmatrix}
\]

    \item A real eigenvector \(\mathbf{v}\) satisfies \(\mathbf{M}\mathbf{v}=k\mathbf{v}\) for some real \(k\), so it is mapped to a vector parallel to itself. A rotation through an acute angle (so not \(0^\circ\) or \(180^\circ\)) turns every vector not parallel to the axis into a vector in a different direction, not parallel to the original. So the only vectors mapped to multiples of themselves lie along the axis, which means the only real eigenvectors are non-zero multiples of \(\mathbf{e}\).
  \end{enumerate}
\end{enumerate}
\end{djsSolution}

\newpage

% ---- Question 13 ----
\questionitem
% [Source: _QBT___Solns__Eigenvalues__Eigenvectors_and_Diagonalisation, Q13]

The matrix $\mathbf{A}$ is given by
\[\mathbf{A} = \begin{pmatrix} 6 & -2 & 1 \\ -2 & 3 & 2 \\ 1 & 2 & 6 \end{pmatrix}\]
\begin{enumerate}[label=\textbf{(\alph*)}]
  \item Find the characteristic equation of $\mathbf{A}$. Hence show that $1$ is an eigenvalue of $\mathbf{A}$, and that $\mathbf{A}$ has one other eigenvalue, which is repeated. State the value of this repeated eigenvalue. \marks{5} \\
  \item Find an eigenvector of $\mathbf{A}$ corresponding to the eigenvalue $1$. Find also two eigenvectors of $\mathbf{A}$, corresponding to the repeated eigenvalue, which are perpendicular to each other. \marks{6} \\
  \item Hence find an orthogonal matrix $\mathbf{P}$ and a diagonal matrix $\mathbf{D}$ such that $\mathbf{P}^{\mathrm{T}}\mathbf{A}\mathbf{P} = \mathbf{D}$. \marks{3} \\
\end{enumerate}

\vspace*{10pt}

\begin{djsSolution}
\begin{enumerate}[label=\textbf{(\alph*)}, start=1]
  \item The characteristic equation is \(\det(\mathbf{A}-\lambda\mathbf{I})=0\), where
\[
\det(\mathbf{A}-\lambda\mathbf{I})=\begin{vmatrix} 6-\lambda & -2 & 1 \\ -2 & 3-\lambda & 2 \\ 1 & 2 & 6-\lambda \end{vmatrix}
\]

Expanding along the first row:
\begin{align*}
\det(\mathbf{A}-\lambda\mathbf{I}) &= (6-\lambda)\big[(3-\lambda)(6-\lambda)-4\big]+2\big[-2(6-\lambda)-2\big]+1\big[-4-(3-\lambda)\big]\\
&= (6-\lambda)(\lambda^2-9\lambda+14)+2(2\lambda-14)+(\lambda-7)\\
&= -\lambda^3+15\lambda^2-68\lambda+84+4\lambda-28+\lambda-7\\
&= -\lambda^3+15\lambda^2-63\lambda+49
\end{align*}

So the characteristic equation is
\[
\lambda^3-15\lambda^2+63\lambda-49=0
\]

Substituting \(\lambda=1\):
\[
1-15+63-49=0
\]

so \(1\) is an eigenvalue of \(\mathbf{A}\), and \((\lambda-1)\) is a factor. Factorising:
\[
\lambda^3-15\lambda^2+63\lambda-49=(\lambda-1)(\lambda^2-14\lambda+49)=(\lambda-1)(\lambda-7)^2
\]

Hence the only other eigenvalue is \(\lambda=7\), which is repeated.

  \item For \(\lambda=1\), solve \((\mathbf{A}-\mathbf{I})\begin{pmatrix}x\\y\\z\end{pmatrix}=\mathbf{0}\), where
\[
\mathbf{A}-\mathbf{I}=\begin{pmatrix} 5 & -2 & 1 \\ -2 & 2 & 2 \\ 1 & 2 & 5 \end{pmatrix}
\]

The first two rows give \(5x-2y+z=0\) and \(-x+y+z=0\). From the second, \(z=x-y\), so \(5x-2y+x-y=0\), giving \(y=2x\) and \(z=-x\). This also satisfies the third row, \(x+4x-5x=0\). Taking \(x=1\), an eigenvector is
\[
\begin{pmatrix}1\\2\\-1\end{pmatrix}
\]

For \(\lambda=7\):
\[
\mathbf{A}-7\mathbf{I}=\begin{pmatrix} -1 & -2 & 1 \\ -2 & -4 & 2 \\ 1 & 2 & -1 \end{pmatrix}
\]

Every row is a multiple of \(x+2y-z\), so the eigenvectors are the non-zero vectors satisfying \(x+2y-z=0\). One such vector is \(\begin{pmatrix}1\\0\\1\end{pmatrix}\), since \(1+0-1=0\).

For a second eigenvector \(\begin{pmatrix}p\\q\\r\end{pmatrix}\) perpendicular to this one, we need
\[
p+2q-r=0,\qquad p+r=0
\]

So \(r=-p\) and \(2p+2q=0\), giving \(q=-p\). Taking \(p=1\), two perpendicular eigenvectors for \(\lambda=7\) are
\[
\begin{pmatrix}1\\0\\1\end{pmatrix},\qquad \begin{pmatrix}1\\-1\\-1\end{pmatrix}
\]

  \item The three eigenvectors from part (b) are mutually perpendicular: the two for \(\lambda=7\) were chosen perpendicular, and each satisfies \(x+2y-z=0\), so has zero scalar product with \(\begin{pmatrix}1\\2\\-1\end{pmatrix}\). Their magnitudes are \(\sqrt{6}\), \(\sqrt{2}\) and \(\sqrt{3}\) respectively. Using the unit eigenvectors as columns gives an orthogonal matrix, so \(\mathbf{P}^{-1}=\mathbf{P}^{\mathrm{T}}\) and \(\mathbf{P}^{\mathrm{T}}\mathbf{A}\mathbf{P}=\mathbf{D}\), with the eigenvalues in \(\mathbf{D}\) in the same order as the columns of \(\mathbf{P}\):
\[
\mathbf{P}=\begin{pmatrix} \dfrac{1}{\sqrt{6}} & \dfrac{1}{\sqrt{2}} & \dfrac{1}{\sqrt{3}} \\[6pt] \dfrac{2}{\sqrt{6}} & 0 & -\dfrac{1}{\sqrt{3}} \\[6pt] -\dfrac{1}{\sqrt{6}} & \dfrac{1}{\sqrt{2}} & -\dfrac{1}{\sqrt{3}} \end{pmatrix},\qquad \mathbf{D}=\begin{pmatrix} 1 & 0 & 0 \\ 0 & 7 & 0 \\ 0 & 0 & 7 \end{pmatrix}
\]

\end{enumerate}
\end{djsSolution}

\newpage

% ---- Question 14 ----
\questionitem
% [Source: _QBT___Solns__Eigenvalues__Eigenvectors_and_Diagonalisation, Q14]

\[\mathbf{A} = \begin{pmatrix} -3 & 2 & -2 \\ -2 & 2 & p \\ 10 & -5 & 3 \end{pmatrix} \quad \text{where } p \text{ is a constant}\] \\
Given that $\begin{pmatrix} 1 \\ 2 \\ -1 \end{pmatrix}$ is an eigenvector for $\mathbf{A}$, \\
\begin{enumerate}[label=\textbf{(\alph*)}]
  \item
  \begin{enumerate}[label=(\roman*)]
    \item determine the eigenvalue corresponding to this eigenvector, \marks{1}
    \item hence show that $p = -4$, \marks{2}
    \item determine the remaining eigenvalues and corresponding eigenvectors of $\mathbf{A}$. \marks{7} \\
  \end{enumerate}
  \item Write down a matrix $\mathbf{P}$ and a diagonal matrix $\mathbf{D}$ such that $\mathbf{P}^{-1}\mathbf{A}\mathbf{P} = \mathbf{D}$. \marks{1} \\
  \item
  \begin{enumerate}[label=(\roman*)]
    \item Solve the differential equation $\dot{u} = ku$, where $k$ is a constant. \marks{2} \\
    With respect to a fixed origin $O$, the velocity of a particle moving through space is modelled by
    \[\begin{pmatrix} \dot{x} \\ \dot{y} \\ \dot{z} \end{pmatrix} = \mathbf{A} \begin{pmatrix} x \\ y \\ z \end{pmatrix}\]
    By considering
    \[\begin{pmatrix} u \\ v \\ w \end{pmatrix} = \mathbf{P}^{-1} \begin{pmatrix} x \\ y \\ z \end{pmatrix} \quad \text{so that} \quad \begin{pmatrix} \dot{u} \\ \dot{v} \\ \dot{w} \end{pmatrix} = \mathbf{P}^{-1} \begin{pmatrix} \dot{x} \\ \dot{y} \\ \dot{z} \end{pmatrix}\]
    \item determine a general solution for the displacement of the particle. \marks{4} \\
    At time $t = 0$ the particle is at the point with position vector $\begin{pmatrix} 3 \\ 7 \\ -1 \end{pmatrix}$. \\
    \item Find expressions for $x$, $y$ and $z$ in terms of $t$ for this particle. \marks{2} \\
  \end{enumerate}
\end{enumerate}

\vspace*{10pt}

\begin{djsSolution}
\begin{enumerate}[label=\textbf{(\alph*)}, start=1]
  \item
  \begin{enumerate}[label=(\roman*), start=1]
    \item
\[
\mathbf{A}\begin{pmatrix}1\\ 2\\ -1\end{pmatrix}=\begin{pmatrix}-3+4+2\\ -2+4-p\\ 10-10-3\end{pmatrix}=\begin{pmatrix}3\\ 2-p\\ -3\end{pmatrix}
\]

The first and third components are \(3\times 1\) and \(3\times(-1)\), so the eigenvalue is \(3\).

    \item Since the vector is an eigenvector with eigenvalue \(3\), the second component must equal \(3\times 2\):
\[
2-p=6\quad\Longrightarrow\quad p=-4
\]

as required.

    \item With \(p=-4\), the matrix is
\[
\mathbf{A}=\begin{pmatrix}-3&2&-2\\ -2&2&-4\\ 10&-5&3\end{pmatrix}
\]

The eigenvalues satisfy \(\det(\mathbf{A}-\lambda\mathbf{I})=0\). Expanding along the first row:
\begin{align*}
\det(\mathbf{A}-\lambda\mathbf{I})&=(-3-\lambda)\big[(2-\lambda)(3-\lambda)-20\big]-2\big[-2(3-\lambda)+40\big]-2\big[10-10(2-\lambda)\big]\\
&=-(\lambda+3)(\lambda^2-5\lambda-14)-2(2\lambda+34)-2(10\lambda-10)\\
&=-\lambda^3+2\lambda^2+29\lambda+42-4\lambda-68-20\lambda+20\\
&=-\lambda^3+2\lambda^2+5\lambda-6
\end{align*}

So the characteristic equation is
\[
\lambda^3-2\lambda^2-5\lambda+6=0
\]

Since \(\lambda=3\) is a root, \((\lambda-3)\) is a factor:
\[
(\lambda-3)(\lambda^2+\lambda-2)=0\quad\Longrightarrow\quad(\lambda-3)(\lambda-1)(\lambda+2)=0
\]

so the remaining eigenvalues are \(\lambda=1\) and \(\lambda=-2\).

\textbf{\(\lambda=1\):} solve \((\mathbf{A}-\mathbf{I})\begin{pmatrix}\alpha\\ \beta\\ \gamma\end{pmatrix}=\mathbf{0}\), that is
\[
\begin{pmatrix}-4&2&-2\\ -2&1&-4\\ 10&-5&2\end{pmatrix}\begin{pmatrix}\alpha\\ \beta\\ \gamma\end{pmatrix}=\begin{pmatrix}0\\ 0\\ 0\end{pmatrix}
\]

The first two rows give
\begin{align*}
-2\alpha+\beta-\gamma&=0\\
-2\alpha+\beta-4\gamma&=0
\end{align*}

Subtracting gives \(3\gamma=0\), so \(\gamma=0\) and \(\beta=2\alpha\). The third row, \(10\alpha-5\beta+2\gamma=10\alpha-10\alpha=0\), is then satisfied. Taking \(\alpha=1\), an eigenvector is
\[
\begin{pmatrix}1\\ 2\\ 0\end{pmatrix}
\]

\textbf{\(\lambda=-2\):} solve \((\mathbf{A}+2\mathbf{I})\begin{pmatrix}\alpha\\ \beta\\ \gamma\end{pmatrix}=\mathbf{0}\), that is
\[
\begin{pmatrix}-1&2&-2\\ -2&4&-4\\ 10&-5&5\end{pmatrix}\begin{pmatrix}\alpha\\ \beta\\ \gamma\end{pmatrix}=\begin{pmatrix}0\\ 0\\ 0\end{pmatrix}
\]

The first and third rows give
\begin{align*}
-\alpha+2\beta-2\gamma&=0\\
2\alpha-\beta+\gamma&=0
\end{align*}

From the first, \(\alpha=2\beta-2\gamma\). Substituting into the second gives \(4\beta-4\gamma-\beta+\gamma=3\beta-3\gamma=0\), so \(\beta=\gamma\) and \(\alpha=0\). The second row is a multiple of the first, so is also satisfied. Taking \(\beta=1\), an eigenvector is
\[
\begin{pmatrix}0\\ 1\\ 1\end{pmatrix}
\]

  \end{enumerate}
  \item Using the eigenvectors as columns, in the same order as their eigenvalues on the diagonal:
\[
\mathbf{P}=\begin{pmatrix}1&1&0\\ 2&2&1\\ -1&0&1\end{pmatrix},\qquad \mathbf{D}=\begin{pmatrix}3&0&0\\ 0&1&0\\ 0&0&-2\end{pmatrix}
\]

  \item
  \begin{enumerate}[label=(\roman*), start=1]
    \item Separating the variables:
\begin{align*}
\int\frac{1}{u}\,\mathrm{d}u&=\int k\,\mathrm{d}t\\
\ln|u|&=kt+c\\
u&=Ke^{kt}
\end{align*}

where \(K\) is an arbitrary constant.

    \item Using the given relations and \(\begin{pmatrix}x\\ y\\ z\end{pmatrix}=\mathbf{P}\begin{pmatrix}u\\ v\\ w\end{pmatrix}\):
\begin{align*}
\begin{pmatrix}\dot{u}\\ \dot{v}\\ \dot{w}\end{pmatrix}&=\mathbf{P}^{-1}\begin{pmatrix}\dot{x}\\ \dot{y}\\ \dot{z}\end{pmatrix}=\mathbf{P}^{-1}\mathbf{A}\begin{pmatrix}x\\ y\\ z\end{pmatrix}\\
&=\mathbf{P}^{-1}\mathbf{A}\mathbf{P}\begin{pmatrix}u\\ v\\ w\end{pmatrix}=\mathbf{D}\begin{pmatrix}u\\ v\\ w\end{pmatrix}
\end{align*}

So the equations separate as
\[
\dot{u}=3u,\qquad \dot{v}=v,\qquad \dot{w}=-2w
\]

By part (c)(i), we have
\[
u=C_1e^{3t},\qquad v=C_2e^{t},\qquad w=C_3e^{-2t}
\]

where \(C_1\), \(C_2\), \(C_3\) are arbitrary constants. Therefore
\begin{align*}
\begin{pmatrix}x\\ y\\ z\end{pmatrix}&=\begin{pmatrix}1&1&0\\ 2&2&1\\ -1&0&1\end{pmatrix}\begin{pmatrix}C_1e^{3t}\\ C_2e^{t}\\ C_3e^{-2t}\end{pmatrix}\\
&=C_1e^{3t}\begin{pmatrix}1\\ 2\\ -1\end{pmatrix}+C_2e^{t}\begin{pmatrix}1\\ 2\\ 0\end{pmatrix}+C_3e^{-2t}\begin{pmatrix}0\\ 1\\ 1\end{pmatrix}
\end{align*}

    \item When \(t=0\):
\begin{align*}
C_1+C_2&=3\\
2C_1+2C_2+C_3&=7\\
-C_1+C_3&=-1
\end{align*}

The first two equations give \(C_3=7-6=1\), then the third gives \(C_1=2\), and so \(C_2=1\). Hence
\[
x=2e^{3t}+e^{t},\qquad y=4e^{3t}+2e^{t}+e^{-2t},\qquad z=-2e^{3t}+e^{-2t}
\]

  \end{enumerate}
\end{enumerate}
\end{djsSolution}

\end{enumerate}
\end{document}
